% SPDX-FileCopyrightText: 2026 Will Cook
% SPDX-License-Identifier: CC-BY-4.0
%
% Long-form reasoning record seeded from the complete problem note.
% The short note and this record are now independent publication units.
% This flat manuscript is generated from reviewable parts below.
\documentclass[11pt]{article}
% SPDX-FileCopyrightText: 2026 Will Cook
% SPDX-License-Identifier: CC-BY-4.0
%
% Shared preamble for the Erdős Problem Notes: the short, standalone,
% problem-owned manuscripts covering the expansion library `ErdosProblems`.
%
% One note owns one Erdős problem.  Each note repeats whatever context its own
% reader needs, so that a reader who arrives at a single problem never has to
% assemble it from the gateway paper.  Deliberate repetition across notes is the
% design, not an oversight.
%
% Source links resolve against one pinned formal-source commit, declared once
% below.  `scripts/check_problem_note_sources.py` verifies that every authored
% (file, line, declaration) triple names that exact declaration in the pinned
% snapshot, so a link cannot silently rot as later work moves lines.

\usepackage{paper-house-style}
\usepackage{array}
\usepackage{booktabs}
\usepackage{enumitem}
\setlist{topsep=0.35em,itemsep=0.2em}
\usepackage[colorlinks=true,linkcolor=linkinternal,urlcolor=linkexternal,
  citecolor=linkinternal]{hyperref}
\hypersetup{bookmarksnumbered=true,bookmarksopen=true,bookmarksopenlevel=1,
  pdfauthor={Will Cook},pdflang={en-GB}}

% ---------------------------------------------------------------------------
% Pinned formal source.  \commit names the checkpoint every link resolves
% against; the checker reads that snapshot rather than the working tree, so the
% printed coordinates stay correct however later waves move declarations.
% ---------------------------------------------------------------------------
\newcommand{\commit}{e5fd26a6d1b1a346430205eab5385b4c9565d004}
\newcommand{\commitshort}{e5fd26a6d1b1}
\newcommand{\repobase}{https://github.com/wcook04/plectis-lean-erdos249-257/blob/\commit}
\newcommand{\sourceurl}{https://github.com/wcook04/plectis-lean-erdos249-257/tree/\commit}
\newcommand{\PX}{ErdosProblems}

% Declaration names stay inside link targets.  The printed label is either a
% spaced, lower-case reading of the declaration name (\lref) or a short authored
% phrase (\lword); neither prints a module path or a raw Lean identifier.
\ExplSyntaxOn
\cs_new_protected:Npn \noteleanlabel #1
  {
    \tl_set:Nx \l_tmpa_tl { \tl_to_str:n {#1} }
    \regex_replace_all:nnN { _+ } { \c{space} } \l_tmpa_tl
    \regex_replace_all:nnN
      { ([a-z0-9])([A-Z]) }
      { \1\c{space}\2 }
      \l_tmpa_tl
    \tl_set:Nx \l_tmpa_tl { \tl_lower_case:n { \tl_use:N \l_tmpa_tl } }
    \tl_use:N \l_tmpa_tl
  }
\ExplSyntaxOff
\newcommand{\leanlabel}[1]{\noteleanlabel{#1}}
\newcommand{\lref}[3]{\href{\repobase/\PX/#1\#L#2}{\leanlabel{#3}}}
\newcommand{\lword}[4]{\href{\repobase/\PX/#1\#L#2}{#4}}
\newcommand{\lloc}[2]{\href{\repobase/\PX/#1\#L#2}{Lean source}}

% Links into a sibling library, named repository-relative.  The expansion
% library is implicit in \lref/\lword, because that is where problem-owned
% work lands.  Several problems have their headline theorem in the older
% reviewed #249/#257 corpus, and a note that could not name that library
% would have to paraphrase a checked statement instead of linking it.
\newcommand{\mref}[3]{\href{\repobase/#1\#L#2}{\leanlabel{#3}}}
\newcommand{\mword}[4]{\href{\repobase/#1\#L#2}{#4}}
\newcommand{\mloc}[2]{\href{\repobase/#1\#L#2}{Lean source}}
\emergencystretch=2em

\newtheorem{theorem}{Theorem}[section]
\newtheorem{proposition}[theorem]{Proposition}
\newtheorem{lemma}[theorem]{Lemma}
\newtheorem{corollary}[theorem]{Corollary}
\theoremstyle{definition}
\newtheorem{definition}[theorem]{Definition}
\newtheorem{example}[theorem]{Example}
\newtheorem{problem}[theorem]{Problem}
\theoremstyle{remark}
\newtheorem*{remark}{Remark}

\newcommand{\N}{\mathbb{N}}
\newcommand{\Npos}{\mathbb{N}_{>0}}
\newcommand{\Z}{\mathbb{Z}}
\newcommand{\Q}{\mathbb{Q}}
\newcommand{\R}{\mathbb{R}}
\newcommand{\lcm}{\operatorname{lcm}}
\newcommand{\ph}{\varphi}

% The series line printed under every note's title block.
\newcommand{\seriesline}{Erd\H{o}s Problem Notes}

% Production provenance.  The wording is fixed by the corpus provenance
% contract and reproduced verbatim in every manuscript in this repository, so
% the papers cannot drift into different accounts of how they were made.
\newcommand{\noteprovenance}{\provenancenote{\emph{Authorship and AI use.} The
prose, formal proofs, and software in this version were drafted and revised by
large-language-model agents under Will Cook's direction.  Cook set the
objectives and acceptance criteria, reviewed and authorised the manuscript, and
is responsible for its contents.  The agents are production tools, not authors.
See \emph{Statements and declarations}.}}

% The status paragraph every note carries, in its own words, before any result.
% Kernel checking establishes the exact Lean proposition; it does not establish
% that the proposition is the interesting one, that it is new, or that the
% Erdős problem is closed.
\newcommand{\statusboundary}{%
  \emph{Status.}  The problem treated here is open, and this note does not
  close it.  Every statement below marked as checked is a proposition that the
  pinned Lean kernel accepts from the sources this note links to, with no
  \texttt{sorry}, no added axiom, and no unchecked evaluation.  That is a claim
  about the formal statement, not about its mathematical interest, its
  novelty, or the original problem.  The unresolved obligations are named
  exactly, in their own section, and none of the finite computations,
  reductions, or no-go results here removes one of them.}

% Problem-specific source pin: #243's cancellation and feedback interfaces
% advance independently of the corpus default.
\renewcommand{\commit}{4ab50e144847a3393fb340341760e5d6907d702e}
\renewcommand{\commitshort}{4ab50e144847}
\hypersetup{
  pdftitle={Reciprocal-tail rigidity: complete reasoning record for Erdős Problem 243},
  pdfsubject={Erdős Problem Notes: Problem 243},
  pdfkeywords={irrationality, Ahmes series, Sylvester sequence, unit fractions, Lean 4},
}

% Names for the state maps and the two search objects.  They are set as
% upright operators so that they read as maps rather than as products.
\newcommand{\sylnext}{\operatorname{syl}}
\newcommand{\den}{\operatorname{den}}
\newcommand{\tail}{\operatorname{tail}}
\newcommand{\ctr}{\operatorname{ctr}}
\newcommand{\defect}{\operatorname{def}}
\newcommand{\fnum}{\operatorname{num}}
\newcommand{\surv}{\operatorname{Surv}}

\runninghead{Erd\H{o}s \#243: complete reasoning record}
\title{Reciprocal-Tail Rigidity: Complete Reasoning Record}
\subtitle{Erd\H{o}s Problem~\#243}
\author{Will Cook}
\date{July 2026}

\begin{document}
% ---- part core ----
\maketitle
\noteprovenance

\begin{abstract}
\noindent
Sylvester's sequence satisfies $a_{n+1}=a_n^2-a_n+1$ and has reciprocal sum
$1$.  Erd\H{o}s asked whether every increasing integer sequence with
$a_{n+1}/a_n^2\to1$ and rational reciprocal sum satisfies this recurrence
eventually.  The leading result of this note carries no rationality hypothesis
and no state hypothesis: if strictly increasing positive integers satisfy
$a_n^2/a_{n+1}=1+3/n+o(n^{-3})$, then $\sum_n1/a_n$ is irrational
(Theorem~\ref{res:cubicrate}).  That rate lies inside the growth hypothesis of
Problem~\ref{res:problem}, since $1+3/n+o(n^{-3})\to1$, so on that rate class
the assertion of the problem holds and no counterexample remains to be checked;
the conclusion is reached by excluding a rational reciprocal sum rather than by
producing the recurrence, and Sylvester's own sequence is not in the class.
Every other rate remains undecided, and Problem~\#243 itself remains open.  The
proof is ordinary mathematics resting on a rising-factorial cubic exclusion
(Theorem~\ref{res:cubicexclusion}), with an exact finite replay of every finite
claim, and no Lean declaration covers either statement.

The strongest statement carrying the hypotheses of Problem~\ref{res:problem}
and no others is a target-equivalent reformulation.  Write $(C_n,D_n,E_n)$ for the integer state
of Section~\ref{sec:state}, $L_n=\lcm(q,a_0,\ldots,a_{n-1})$, $M_n=D_n/L_n$,
$U_n=C_n/M_n$, $V_n=E_n/M_n$, and $\mathcal R$ for the set of steps that set a
new record for $U$, where $q$ is a common denominator of the reciprocal sum.
Then the sequence is eventually Sylvester if and only if
$\sum_{n\in\mathcal R}(-V_n-B)_+f(U_n)$ is finite for some integer $B\ge0$,
where $f$ is any finite nonincreasing weight with divergent integral
(Theorem~\ref{res:weightedrecord}).  Equation~\eqref{eq:weightedgrowth} is an
equivalent form written in the growth defect $a_n^2/a_{n+1}-1$.

That criterion is a target-equivalent reformulation, and it is an ordinary
mathematical proof.  Its finite first-crossing
arithmetic is kernel-checked, in the
\lword{Erdos243/LcmRecordExcess.lean}{31}{covered_wall_forces_excess}{covered-wall
charge} and the three lemmas beside it, and the analytic half is not
formalised.  The criterion settles nothing on its own, because the growth
hypothesis does not supply the finiteness, and termwise convergence to zero is
insufficient.

Section~\ref{sec:records} carries further criteria in the record coordinates,
each with its hypotheses and its evidence class.  Write
$\ell(x)=\log_2\log_2\max(4,x)$ and $H_n=\max_{j\le n}C_j$.  Two are exact
equivalences for the endpoint under arbitrary cancellation, one through the
record-amplified negative error and one through unit true record increments.
One is a dichotomy for the record-increment coefficient
$\Theta=\limsup_n(H_{n+1}-H_n)/\ell(H_n)$, which is $0$ or strictly greater
than $1$, so $\Theta\le1$ already forces the recurrence.  One replaces a constant bound
on the negative part by the bound $(1-\delta)\ell(C_n)$.  One states the
Lean-checked endpoint in the language of the original problem, where the
Erd\H{o}s--Straus quantity $Q_n$ is asked only for $\limsup_nQ_n<\infty$ in
place of their $\limsup_nQ_n\le0$.  One is the cubic-rate irrationality
theorem stated above, with the rising-factorial exclusion it rests on.  The
last two
record what the method cannot reach: the classical Ahmes criteria all decide
exactly the half-space $E_n\ge0$ eventually, and pairwise coprimality with
whole-modulus avoidance admits rises of order $\log C_n$.

Two conditional endpoints on the state system are proved in Lean.  Under the
exact $C/D$ dynamics, positivity, normalised vanishing, and an eventually
bounded negative part, $E_n=0$ eventually, so the Sylvester recurrence holds
eventually
\lword{Erdos243/ReciprocalTailRigidity.lean}{2403}
      {boundedNegativePart_sylvesterNext_eventually}
      {the bounded-negative-part theorem}.  Normalised vanishing itself yields
the required strict centring after a finite shift.  The same argument in
summable form gives a second criterion: under the exact dynamics, positivity,
the strict centred-step equation, normalised vanishing, and
$\sum_n(-E_n)_+/C_n<\infty$, the Sylvester recurrence is recovered
\lword{Erdos243/SparseResetRecovery.lean}{510}
      {sylvesterNext_eventually_of_summable_negativeRelativeMass}
      {the finite-negative-mass recovery theorem}.
Neither hypothesis is known for the canonical orbit.  The proof of the first
has three arithmetic steps.  Bounded negative errors give bounded upward
increments of $C_n$.  On a hypothetical tail with no zero error, integrality
and normalised vanishing force $C_n\to\infty$; the tail gcd then stabilises.
After reduction, the successive multipliers are pairwise coprime and each is
coprime to all later numerators; a block of consecutive multiples contradicts
these coprimality conditions.  Finally, $E_n=0$ is absorbing.

Several sharper exclusions explain why the remaining case is genuinely
unbounded.  A negative error cannot be constant, at any magnitude or scale, and
cannot be periodic when $|E_n|<a_n$ and the drift is positive.  If $E_n\ge0$,
the state is nonincreasing and therefore eventually zero
\lword{Erdos243/ReciprocalTailRigidity.lean}{1837}{centeredState_eventually_zero}{nonnegative-error descent};
the recurrence then follows from
\lword{Erdos243/ReciprocalTailRigidity.lean}{1799}{sylvesterNext_eventually_of_centered_zero}{eventual Sylvester recurrence}.
On a reduced recovery interval of length $L$, deleted old moduli satisfy the
exact inequality
\[
 K^L\lcm(m_q:q\in R)^2<(K+1)^L.
\]
Hence fixed-length recoveries are eventually cancellation-free, although
recoveries of increasing length remain possible.

One clause of that frontier is already published and is credited as such.  That
a counterexample must have negative errors infinitely often is the
contrapositive of Koizumi's Proposition~19(2)~\cite[Prop.~19(2)]{koizumi2025},
which he attributes to Badea; Koizumi and Badea assume the sign is eventually
nonnegative.  The theorem proved here assumes only that the negative part is
eventually bounded, which admits sign changes of bounded depth at every scale,
so it is a strengthening of their hypothesis rather than a new clause.  The
unboundedness and divergent-mass parts of the frontier have no located
antecedent.

Koizumi's pseudo-greedy coordinates supply normalised vanishing and a stronger
centring range for the canonical Erd\H{o}s orbit, but neither boundedness nor
summability of its negative part.  Thus any surviving orbit must have
cofinally unbounded negative excursions and divergent normalised negative
mass.  Excluding that case is the exact open boundary, and Problem~\#243
remains open.
\end{abstract}

\section{Introduction}\label{sec:problem}

Sylvester's sequence $2,3,7,43,1807,\ldots$ is generated by
$a_{n}=a_{n-1}^{2}-a_{n-1}+1$ and has reciprocal sum $1$; the recurrence is
exactly the statement that the tail beyond $a_{n-1}$ equals $1/(a_{n-1}-1)$.
Erd\H{o}s Problem~\#243 asserts that this is the only way a sequence of that
growth can have a rational reciprocal sum:

\begin{problem}[Erd\H{o}s \#243]\label{res:problem}
Let $1\le a_1<a_2<\cdots$ be a sequence of integers with
\[
 \lim_{n\to\infty}\frac{a_n}{a_{n-1}^{2}}=1
 \qquad\text{and}\qquad
 \sum\frac{1}{a_n}\in\Q .
\]
Then $a_n=a_{n-1}^{2}-a_{n-1}+1$ for all sufficiently large $n$.
\end{problem}

\noindent See~\cite[p.~64]{erdosgraham1980} and~\cite[p.~105]{erdos1988}.
Bloom's catalogue records the problem as open~\cite{erdosproblems}.
Below we index the same recurrence as
$a_{n+1}=a_n^{2}-a_n+1$, which is the shift the formal sources use; the two
forms are the same statement.  Write $\sylnext(a)=a^{2}-a+1$, the
\lword{Erdos243/ReciprocalTailRigidity.lean}{37}{sylvesterNext}{Sylvester
successor}.

\paragraph{Relation to prior work.}
Two published results reach the original problem under analytic hypotheses that
are not assumed below.  Their hypotheses are conditions on the sequence $(a_n)$
itself.  Duverney's is not comparable with the state-level hypotheses used
here; the Erd\H{o}s--Straus one, read through the dictionary of the next
paragraph, is implied by nonnegativity of $E$, the hypothesis of
Theorem~\ref{res:descent}, so that criterion already covers the descent proved
below (see~\cite[Cor.~20(1) and Prop.~19(1), pp.~12--13]{koizumi2025}).
Erd\H{o}s and Straus proved that if $\lim a_n/a_{n-1}^{2}=1$, the reciprocal
sum is rational, and $a_n$ does not satisfy the recurrence, then
\[
 \limsup_{n\to\infty}\ \frac{[a_1,\ldots,a_n]}{a_{n+1}}
 \left(\frac{a_{n+1}^{2}}{a_{n+2}}-1\right)>0,
\]
where $[a_1,\ldots,a_n]$ is the least common
multiple~\cite[Theorem~3, p.~132]{erdosstraus1964}.  This is the indexing in
the original theorem.  Erd\H{o}s's 1988 survey, followed by the current
catalogue summary, instead prints $a_n^2/a_{n+1}-1$ in the second factor while
retaining the same prefix-LCM quotient; that displayed summary is off by one
and is not followed here~\cite[p.~105]{erdos1988}\cite{erdosproblems}.
Koizumi records the convenient sufficient rate
$a_n^2/a_{n+1}=1+o(1/n)$ under which the Erd\H{o}s--Straus criterion settles
the problem~\cite[Remark~21, p.~13]{koizumi2025}.

Duverney proved a
conditional signed form of the problem itself: for signs
$\epsilon_n\in\{-1,1\}$, if
$\sum_{n\ge0}(a_{n+1}/a_n^{2}-1)$ converges, then the reciprocal sum with
numerators $\epsilon_n$ is rational if and only if
\[
 a_{n+1}=a_n^2-(\epsilon_{n+1}/\epsilon_n)a_n
             +\epsilon_{n+2}/\epsilon_{n+1}
\]
for all large $n$~\cite[Corollary~3.2, p.~287]{duverney2001}.  Absolute
convergence is a stronger sufficient specialisation, not the hypothesis
printed in that corollary.  The all-positive specialisation is the form
relevant here.

Badea's positive-term criterion is adjacent but different:
for positive rational terms $b_n/a_n$, eventual strict inequality
\[
 a_{n+1}>\frac{b_{n+1}}{b_n}a_n^2-\frac{b_{n+1}}{b_n}a_n+1
\]
forces irrationality, while rationality under the corresponding non-strict
inequality forces eventual equality~\cite[Theorem~A, p.~313; Cor.~2.2,
p.~316]{badea1993}.  These positive-term hypotheses do not cover the
mixed-sign state orbits studied here.  Neither result is formalised here, and
the results below do not recover either of them.

Koizumi's work on doubly exponential sequences~\cite{koizumi2025} bears
directly on the state system used below, and supplies the coordinates in which
several of the statements of this note are most naturally read.  For a rational
$r=p/q$ with pseudo-greedy expansion $(a_n)$, remainders $(x_n)$ and gap
sequence $\varepsilon_n=x_n^{-1}+1-a_n$, his Lemma~15
~\cite[p.~9]{koizumi2025} produces integers $c_n>0$,
$d_n$ and $e_n$ with $x_n=c_n/d_n$ and $\varepsilon_n=e_n/c_n$, satisfying
\[
 a_n=\frac{d_n-e_n}{c_n}+1,\qquad
 c_{n+1}=c_n-e_n,\qquad
 d_{n+1}=a_nd_n ,
\]
and the proof of that lemma also gives $c_{n+1}=a_nc_n-d_n$.  Under the
dictionary $(C_n,D_n,E_n)=(c_n,d_n,e_n)$ these are exactly the objects of
Section~\ref{sec:state}: the first relation is the centring map $\ctr$ rewritten
as $e_n=d_n-(a_n-1)c_n$, the second is Proposition~\ref{res:update}, the third
is the denominator update, and the fourth is the tail update.  The lower-case
$e_n$ of Section~\ref{sec:state} is $-E_n$, so it is the negative of Koizumi's
$e_n$ at an index where the centred state is negative.

His Theorem~16
~\cite[p.~11]{koizumi2025} shows
that his Conjecture~6, that for a positive rational $r$ whose gap sequence
satisfies $\varepsilon_n\to0$ one has $\varepsilon_n=0$ for all large $n$,
is equivalent to an affirmative answer to the question of Erd\H{o}s and Graham
recorded as his Question~5, which is Problem~\ref{res:problem} above.
Two of the implications proved below have prior art in these canonical
coordinates: his Lemma~13, that $\varepsilon_n=0$ forces $\varepsilon_{n+1}=0$,
is the absorption of Theorem~\ref{res:absorb}, and his Proposition~19(2), that
$\varepsilon_n\ge0$ for all large $n$ forces $\varepsilon_n=0$ for all large
$n$, is the descent of Theorem~\ref{res:descent}.  These appear in canonical
coordinates in~\cite[Lemma~13, p.~9; Prop.~19(2), p.~12]{koizumi2025}.

The growth hypothesis in Problem~\ref{res:problem} is calibrated by two facts
about Sylvester's sequence.  It satisfies $a_n\approx c_0^{2^{n}}$ with
$c_0=1.2640847\ldots$, and shifting it further produces sequences with
$a_n\approx C^{2^{n}}$ for arbitrarily large $C$ whose reciprocals still sum
to a rational number~\cite[arXiv v4, p.~2]{kovactao2024}.  The classical sufficient
condition for irrationality, $\lim_n a_n^{1/2^{n}}=\infty$, is therefore
sharp.  Kova\v{c} and Tao identify that condition as folklore
~\cite[arXiv v4, p.~2]{kovactao2024}; they attribute the sharpness observation to
Erd\H{o}s (1975).  A sequence with
$a_n/a_{n-1}^{2}\to1$ has $a_n^{1/2^{n}}$
convergent, so that criterion says nothing about the sequences of
Problem~\ref{res:problem}, and rationality is genuinely possible there.
Sylvester's
sequence is \texttt{A000058} in the OEIS and its shifts \texttt{A129871}.  For
the recent literature on irrationality of Ahmes series we refer to Kova\v{c}
and Tao~\cite{kovactao2024}, who resolve several problems of Erd\H{o}s and
Graham drawn from the same two sources cited above, and whose introduction
gives a sample of the intermediate work, including S\'andor (1984) and Badea
(1987).  They do not treat Problem~\#243; the rigidity conclusion asked for
there is not among their results.

The \emph{Formal Conjectures} collection contains a mathematically equivalent
unproved declaration, up to its zero-based indexing
~\cite{formalconjectures243}.  Its summand is $\Q$-valued, so Lean's
\texttt{Summable} hypothesis asserts the existence of a sum in $\Q$; the
finite indexing shift changes that sum only by a rational prefix.  The
declaration therefore does encode the rationality premise, but its proof is
\texttt{sorry}.  Its role here is statement-level prior art; it supplies no
proof authority, and the development below is independent of it.
The machine-checked results concern the state system: they exclude constant and
periodic negative magnitudes, the bounded-rise barrier of
Theorem~\ref{res:barrier}, and the two conditional endpoints
(Theorems~\ref{res:bounded} and~\ref{res:mass}), each stated with the exact
hypotheses that separate it from Problem~\#243.

The hypothesis is asymptotic and the conclusion is exact, so the argument must
convert an analytic rate into an integer obstruction.  The conversion used here
is the classical one: clear denominators along the sequence, so that
rationality makes a tail integral, and then centre that integer at the value it
would take on Sylvester's sequence.  What remains is a single integer error
$E$, and the whole question is whether $E$ can avoid vanishing.

\statusboundary

\begin{quote}\small
\textbf{Strongest result.}
If strictly increasing positive integers satisfy
$a_n^2/a_{n+1}=1+3/n+o(n^{-3})$, then $\sum_n1/a_n$ is irrational
(Theorem~\ref{res:cubicrate}).  The theorem assumes that rate and nothing else.
The rate class sits inside the growth hypothesis of
Problem~\ref{res:problem}, so the assertion of the problem holds on it, by
exclusion of a rational reciprocal sum and not by production of the recurrence.
This is an ordinary proof with an exact finite replay, and no Lean declaration
covers it.

\smallskip
\textbf{Strongest statement under the problem's own hypotheses.}
Under the hypotheses of Problem~\ref{res:problem} and no others, the sequence
is eventually Sylvester if and only if the weighted record-excess sum
$\sum_{n\in\mathcal R}(-V_n-B)_+f(U_n)$ is finite for some integer $B\ge0$
(Theorem~\ref{res:weightedrecord}).  This is an ordinary proof whose finite
first-crossing arithmetic is kernel-checked and whose analytic half is not.
The hypotheses of Problem~\ref{res:problem} do not supply that finiteness, so
the criterion is a target-equivalent reformulation of the problem rather than a
solution of it.

\smallskip
\textbf{Strongest Lean-checked result.}
Under the exact integer dynamics, positivity, normalised vanishing, and an
eventually bounded negative part, the centred error $E_n$ is eventually zero.
The original sequence therefore satisfies the Sylvester recurrence from some
point onward (Theorem~\ref{res:bounded}).  Hypothesis~(5), the bound on the
negative part, is supplied by nothing in this note.

\smallskip
\textbf{Second endpoint.}
The same conclusion holds when the normalised negative mass
$\sum_n(-E_n)_+/C_n$ is finite, provided the strict centred-step equation is
available (Theorem~\ref{res:mass}).

\smallskip
\smallskip
\textbf{What remains open.}
Koizumi's coordinates provide normalised vanishing for the canonical orbit,
but they do not provide either bounded negative excursions or finite
normalised negative mass.  Any surviving counterexample must therefore have
cofinally unbounded negative excursions and divergent normalised negative
mass (Proposition~\ref{res:frontier}), record increments strictly above the
log-log scale (Theorem~\ref{res:recorddichotomy}), and unbounded
record-amplified error (Theorem~\ref{res:recordamplified}).

\smallskip
\textbf{Proof spine.}
The record criterion charges each first crossing of a CRT-selected height to
the record excess at the step that crosses it, and the crossing heights carry
a divergent weight.  For the state-system endpoints, bounded negative errors
control the rise of the integer tail state; a gcd stabilisation reduces the
orbit to pairwise-coprime fresh moduli, and a block of consecutive multiples
then gives the contradiction.  The finite-mass argument replaces the uniform
bound by a summable recovery estimate.
\end{quote}

\noindent The remaining exclusions sharpen this boundary.  By
Theorem~\ref{res:absorb}, a counterexample has $E_n\ne0$ for
every large $n$.  By Theorem~\ref{res:descent} it cannot have $E$ eventually
nonnegative.  If its tail is all negative, Theorems~\ref{res:constant}
and~\ref{res:periodic} exclude, respectively, constant magnitude and periodic
magnitude in the stated regime $e_n<a_n$ with positive drift; they are not a
classification of a mixed-sign tail.  Under normalised vanishing,
Theorems~\ref{res:bounded} and~\ref{res:mass} show that the negative part is
neither bounded nor of finite normalised mass.  What survives is stated in
Section~\ref{sec:open}.

\subsection*{Structure}

Section~\ref{sec:cubicrate} proves the unconditional cubic-rate irrationality
theorem, which is the only result here that settles the assertion of
Problem~\ref{res:problem} on a class of sequences.
Section~\ref{sec:state} sets up the state system, and Section~\ref{sec:defect}
proves the defect identity, from which eventual vanishing of $E$ forces the
Sylvester recurrence and a single vanishing error propagates.
Section~\ref{sec:lcmrecords} then proves the exact record-excess criterion, and
Section~\ref{sec:records} states the further criteria in the record
coordinates.  Section~\ref{sec:descent} disposes of the case $E\ge0$ in three
lines.  The rest of the note is the case $E<0$, in four widening steps:
constant magnitude (Section~\ref{sec:constant}), periodic magnitude
(Section~\ref{sec:periodic}), bounded magnitude via a coprimality barrier
(Sections~\ref{sec:barrier} and~\ref{sec:bounded}), finite normalised negative
mass (Section~\ref{sec:mass}), and what is left (Section~\ref{sec:open}).
These sections contain the new content; a reader who wants only the strongest
statements should read Section~\ref{sec:cubicrate}, then
Sections~\ref{sec:lcmrecords} and~\ref{sec:records}, and
then Section~\ref{sec:open}.  Appendix~\ref{app:index} is a guide to the formal
sources, and Appendix~\ref{app:residue} records an exact residue reduction that
an earlier finite search used.  Linked phrases open the corresponding Lean
declaration at the pinned source revision~\commitshort.  Where a section names
a Lean module that is outside that revision, it says so in place of a link.

\medskip
\noindent\textbf{Keywords.} irrationality; Ahmes series; Sylvester's sequence;
unit fractions; Lean~4.
\qquad
\textbf{MSC 2020.} 11J72 (primary); 11B37, 11D68, 68V20 (secondary).

\section{An unconditional irrationality theorem}\label{sec:cubicrate}

The second theorem of this section is the note's lead, and it is stated in the
sequence alone.  The first is the exclusion it rests on, and it is stated in
the integer state $(C_n,D_n,E_n)$ that Section~\ref{sec:state} constructs from
a rational reciprocal sum; a reader meeting $C_n$ for the first time can read
Theorem~\ref{res:cubicrate} on its own and return here afterwards.

\begin{theorem}[rising-factorial cubic exclusion]\label{res:cubicexclusion}
For every exact positive integer orbit, every $\alpha\in\Q_{>0}$ and every
$\beta\in\Q$,
\[
 \liminf_{X\to\infty}
 \frac{\#\{n\le X:C_n\ne\alpha\,n(n+1)(n+2)+\beta\}}{X}>0 .
\]
\end{theorem}

\begin{theorem}[cubic-rate irrationality]\label{res:cubicrate}
If strictly increasing positive integers satisfy
$a_n^2/a_{n+1}=1+3/n+o(n^{-3})$, then $\sum_n1/a_n$ is irrational.
\end{theorem}

\paragraph{Mechanism.}
Density zero of the exceptions would give four consecutive agreements
arbitrarily late, so the third difference is constant, $G_n$ stabilises, and the
primitive tail carries an integral cubic $Q(n)=(m/6)n(n+1)(n+2)+c$ with
$c=\pm1$ by adjacent coprimality.  For the irreducible case, a numerator zero
makes an explicit quantity a square modulo the specialising prime, so Chebotarev
square specialisation puts $x^2-1$ in the squares of the field generated by a
root $x$ of the associated cubic; the trace conditions then force the sole
integer scale $m=12$.  The two survivors
$2n(n+1)(n+2)\pm1$ have numerator words modulo $7$ that the feedback recursion
cannot realise, which gives lower density at least $1/28$.  For the corollary,
rationality gives the canonical orbit with
$C_{n+1}/C_n=a_n^2/a_{n+1}+O(1/a_n)$ and $1/a_n\ll n^{-3}$; an integer-extraction
lemma at a regular rate then produces an exact rising-factorial cubic, against
the exclusion.

\paragraph{Evidence.}
Ordinary proof with an exact finite replay of every finite claim, including the
two words modulo $7$, the image $\{0,2,5,6\}$ of the recursion, the sole integer
scale $12$, and the countermodel $n(n+1)(n+2)+(-1)^n$, whose ratio is
$1+3/n+O(n^{-3})$ and which is not eventually polynomial.  This countermodel
shows the exponent $o(n^{-\lambda})$ in the extraction lemma is sharp.  No Lean
declaration covers either statement.

\paragraph{Position.}
Theorem~\ref{res:cubicrate} carries no rationality hypothesis and no state
hypothesis, and it is the only unconditional irrationality statement in this
note.  Its rate class lies inside the growth hypothesis of
Problem~\ref{res:problem}, because $a_n^2/a_{n+1}=1+3/n+o(n^{-3})$ forces
$a_{n+1}/a_n^2\to1$.  On that class the assertion of
Problem~\ref{res:problem} therefore holds, because no sequence in it has a
rational reciprocal sum; the recurrence is never exhibited, and Sylvester's own
sequence, whose ratio $a_n^2/a_{n+1}$ is $1+O(1/a_n)$, is not in the class.
The earlier sections decide nothing at this rate, and every other rate inside
the growth hypothesis remains undecided.

\section{The centred integer state}\label{sec:state}

The three maps defined next carry out the two steps announced at the end of
Section~\ref{sec:problem}.  The first two clear denominators along the
sequence, so that a rational reciprocal sum leaves a sequence of integers; the
third subtracts from that integer the value it takes on Sylvester's sequence,
leaving one number whose vanishing is the whole question.  All three are maps
between integer tuples and carry no hypothesis of any kind.

For $a,D,C\in\Z$ define
\[
 \den(a,D)=aD,
 \quad
 \tail(a,D,C)=aC-D,
 \quad
 \ctr(a,D,C)=D-(a-1)C ,
\]
the
\lword{Erdos243/ReciprocalTailRigidity.lean}{41}{nextDenState}{denominator
update}, the
\lword{Erdos243/ReciprocalTailRigidity.lean}{45}{nextTailState}{tail update},
and the
\lword{Erdos243/ReciprocalTailRigidity.lean}{49}{centeredState}{centred state}.
Given a sequence $(a_n)$ of integers, whose terms we call the
\emph{multipliers}, we write $D_{n+1}=\den(a_n,D_n)$ for the \emph{denominator
state}, $C_{n+1}=\tail(a_n,D_n,C_n)$ for the \emph{tail state}, and
$E_n=\ctr(a_n,D_n,C_n)$ for the \emph{centred state}, also called the
\emph{error}.  The word \emph{centred} records what the third map does: it
measures $D_n$ against $(a_n-1)C_n$, and $D_n=(a_n-1)C_n$ holds at every index
on Sylvester's sequence (Example~\ref{ex:sylvester}), so $E$ is a deviation
from that sequence rather than a size.  At an index where $E_n<0$ we write
$e_n=-E_n>0$ and call it the \emph{magnitude} of the error there.

\begin{proposition}[update law]\label{res:update}
$\tail(a,D,C)=C-\ctr(a,D,C)$; that is, $C_{n+1}=C_n-E_n$.
\end{proposition}

\begin{proof}
$C-\bigl(D-(a-1)C\bigr)=aC-D$.
\end{proof}

\noindent Formalised as the
\lword{Erdos243/ReciprocalTailRigidity.lean}{57}{nextTailState_eq_sub_centered}{update
law}.  The state therefore moves by exactly the error at each step, which is
what makes the sign of $E$ decisive: where $E$ is positive the state falls,
where $E$ is negative it rises, and where $E$ vanishes it is stationary.

\begin{example}[the Sylvester orbit]\label{ex:sylvester}
Take $a_n=2,3,7,43,1807,\ldots$ with $a_{n+1}=a_n^{2}-a_n+1$, and start the
state at $D_0=C_0=1$.  The two updates give

\begin{center}
\begin{tabular}{@{}rrrrr@{}}
\toprule
$n$ & $a_n$ & $D_n$ & $C_n$ & $E_n$\\
\midrule
$0$ & $2$ & $1$ & $1$ & $0$\\
$1$ & $3$ & $2$ & $1$ & $0$\\
$2$ & $7$ & $6$ & $1$ & $0$\\
$3$ & $43$ & $42$ & $1$ & $0$\\
$4$ & $1807$ & $1806$ & $1$ & $0$\\
\bottomrule
\end{tabular}
\end{center}

\noindent and $D_n=a_n-1$ with $C_n=1$ at every index: if $D_n=a_n-1$ and
$C_n=1$ then $C_{n+1}=a_n-(a_n-1)=1$ and $D_{n+1}=a_n(a_n-1)=a_{n+1}-1$.
Hence $E_n=D_n-(a_n-1)C_n=0$ throughout and the tail state never moves, which
is Proposition~\ref{res:update} in the stationary case.
\end{example}

\paragraph{The intended reading, which is not formalised.}
Let $T_n=\sum_{k\ge n}1/a_k$ and suppose $T_0\in\Q$.  Put $D_0$ equal to a
common denominator and $D_n=D_0a_0\cdots a_{n-1}$, and set $C_n=D_nT_n$.  Then
$C_n\in\Z$ for every $n$, and from $T_n=1/a_n+T_{n+1}$ one gets
$C_{n+1}=a_nC_n-D_n$, which is the tail update.  On Sylvester's sequence
$T_n=1/(a_n-1)$ exactly, so $D_n=(a_n-1)C_n$ and $E_n=0$: the centred state
measures deviation from the Sylvester tail identity, and it vanishes
identically on the Sylvester orbit.

The Lean module carries no rationality hypothesis and no growth hypothesis; it
proves identities and implications about the integer system above, and about its
natural-number realisation, for which the bridging identities are the
\lword{Erdos243/ReciprocalTailRigidity.lean}{1863}{natTail_eq_nextTailState}{tail
realisation}, the
\lword{Erdos243/ReciprocalTailRigidity.lean}{1875}{natDen_eq_nextDenState}{denominator
realisation}, and the
\lword{Erdos243/ReciprocalTailRigidity.lean}{1972}{natTail_eq_sub_centeredState}{signed
update law}.

The series itself is reached by three further declarations, which were absent
from earlier revisions of this note.  Write $C_n/D_n$ for the real ratio carried
by the natural state, the
\lword{Erdos243/ReciprocalTailRigidity.lean}{1895}{tailRatio}{tail ratio}.
Given $a_n>0$, $D_n>0$, $C_{n+1}+D_n=a_nC_n$ and $D_{n+1}=a_nD_n$, one exact
step removes $1/a_n$ from that ratio
\lword{Erdos243/ReciprocalTailRigidity.lean}{1899}{tailRatio_eq_reciprocal_add_next}{one-step
reciprocal identity}, finite iteration telescopes to a partial sum plus the
endpoint ratio
\lword{Erdos243/ReciprocalTailRigidity.lean}{1923}{tailRatio_eq_partialReciprocalSum_add}{finite
telescoping}, and convergence of the endpoint ratio to zero identifies the
infinite sum
\lword{Erdos243/ReciprocalTailRigidity.lean}{1949}{reciprocalSeries_hasSum_of_tailRatio_tendsto_zero}{the
series realisation}: $\sum_n1/a_n$ then has sum $C_0/D_0$.  These three carry no
growth hypothesis, no centring hypothesis and no rationality hypothesis, and
they exclude nothing.  They close the loop between the integer system and the
series that motivates it.  Producing the state of a given rational sum, which is
the direction the problem needs, is the exposition of this paragraph and is not a
checked statement.

\begin{proposition}[scale equivariance]\label{res:scale}
For every $s,a,D,C\in\Z$,
\[
 \den(a,sD)=s\,\den(a,D),\qquad
 \tail(a,sD,sC)=s\,\tail(a,D,C),
\]
\[
 \ctr(a,sD,sC)=s\,\ctr(a,D,C).
\]
\end{proposition}

\noindent Formalised as the
\lword{Erdos243/ReciprocalTailRigidity.lean}{65}{nextDenState_scale}{denominator
equivariance}, the
\lword{Erdos243/ReciprocalTailRigidity.lean}{70}{nextTailState_scale}{tail
equivariance}, and the
\lword{Erdos243/ReciprocalTailRigidity.lean}{75}{centeredState_scale}{centred
equivariance}.  The system depends only on the ray through $(D,C)$.  This is
used twice below: once to normalise the finite search of
Appendix~\ref{app:residue}, and once, in Theorem~\ref{res:periodic}, as the
induction step that removes common scale.

\section{The defect identity and its two consequences}\label{sec:defect}

Section~\ref{sec:state} connects the Sylvester recurrence to the error in one
direction only: on Sylvester's sequence the error vanishes.  The converse is
what the problem needs, and it comes from a single identity, which expresses
the deviation of $a_{n+1}$
from the Sylvester successor, multiplied by the tail state $C_{n+1}$, as a
combination of the errors at $n$ and at $n+1$, the first weighted by
$a_n^{2}$.  Its two consequences are
Theorem~\ref{res:eventual}, in which an eventually vanishing error forces the
recurrence provided the tail state is eventually nonzero, and
Theorem~\ref{res:absorb}, in which, under strict centring (the condition
$|E_n|<C_n$ at every index), one vanishing error propagates to every later
index.

Write $\defect(a,a')=a'-\sylnext(a)$ for the deviation of the next term from
the Sylvester successor, the
\lword{Erdos243/ReciprocalTailRigidity.lean}{53}{sylvesterDefect}{Sylvester
defect}.

\begin{theorem}[defect identity]\label{res:defect}
For all $a,a',D,C\in\Z$,
\[
 \defect(a,a')\cdot\tail(a,D,C)
 =a^{2}\,\ctr(a,D,C)-\ctr\bigl(a',\den(a,D),\tail(a,D,C)\bigr),
\]
that is, $\Delta_n\,C_{n+1}=a_n^{2}E_n-E_{n+1}$.
\end{theorem}

\begin{proof}
A direct expansion: both sides equal
$a'aC-a'D-a^{3}C+a^{2}D+a^{2}C-aD-aC+D$.
\end{proof}

\noindent Formalised as the
\lword{Erdos243/ReciprocalTailRigidity.lean}{1769}{sylvesterDefect_mul_nextTailState}{defect
identity}.  It is a polynomial identity in four indeterminates, with no
hypothesis of any kind, and it carries the whole argument twice over: once for
rigidity, and once for absorption.

\begin{example}[one defect and the error it creates]\label{ex:defect}
Continue Example~\ref{ex:sylvester} but replace $a_3=43$ by $a_3=44$.  The
states $D_3=42$ and $C_3=1$ are unchanged, since they depend only on
$a_0,a_1,a_2$, and $E_2=0$ still.  The defect is
$\Delta_2=a_3-\sylnext(a_2)=44-43=1$, and
$E_3=D_3-(a_3-1)C_3=42-43=-1$, so the identity reads
\[
 \Delta_2C_3=1\cdot1=1=49\cdot0-(-1)=a_2^{2}E_2-E_3 .
\]
By Proposition~\ref{res:update} the tail state then rises, $C_4=C_3-E_3=2$.  A
single unit of defect at one index has produced a negative error of magnitude
$1$ at the next, and the state has started to climb.
\end{example}

\begin{theorem}[two vanishing errors force the step]\label{res:step}
Let $a,a',D,C\in\Z$ with $\tail(a,D,C)\ne0$.  If
$\ctr(a,D,C)=0$ and $\ctr\bigl(a',\den(a,D),\tail(a,D,C)\bigr)=0$, then
$a'=\sylnext(a)$.
\end{theorem}

\begin{proof}
By Theorem~\ref{res:defect} the product $\defect(a,a')\cdot\tail(a,D,C)$ is
zero, and the second factor is nonzero by hypothesis.
\end{proof}

\noindent Formalised as the
\lword{Erdos243/ReciprocalTailRigidity.lean}{1781}{sylvesterNext_eq_of_centered_zero}{local
rigidity step}.  The hypothesis $C_{n+1}\ne0$ is not removable: at $C_{n+1}=0$
the identity gives no information about $a'$.

\subsection{Two-step cancellation and feedback}

The one-step state equations hide a second-order constraint once the raw
numerator is reduced at two consecutive stages.  Write $u,u_1,u_2$ for the
primitive numerator coordinates, $v,v_1$ for the corresponding denominator
coordinates, and $h,h_1$ for the two common factors removed in reduction.  If
the two steps have multipliers $a,a_1$, then the exact hypotheses are
\[
 hu_1+v=au,\qquad hv_1=av,\qquad
 h_1u_2+v_1=a_1u_1.
\]

\begin{theorem}[two-step primitive recurrence]\label{res:secondorder}
Under these three equations,
\[
 a^2u+hh_1u_2=h(a+a_1)u_1.
\]
\end{theorem}

\begin{proof}
Multiply the first equation by $a$, replace $av$ using the second equation,
and then replace $h_1u_2+v_1$ using the third.  The remaining terms factor as
the displayed right-hand side.
\end{proof}

\noindent This is the
\lword{Erdos243/DynamicCancellation.lean}{270}{reducedStep_secondOrder}{two-step
primitive recurrence}.  The hard point is keeping both changing cancellation
factors visible.  The result is stronger than a one-step congruence, though it
remains non-autonomous and cannot by itself force $a_1=a^2-a+1$.

There is a complementary integer identity when the relevant step is
cancellation-free.  Put $q=ap-p_1$, and suppose the next primitive coordinate
satisfies
\[
 p_2+a^2p=(a+a_1)p_1.
\]
Then
\begin{theorem}[cancellation-free curvature identity]\label{res:curvature}
\[
 q^2+\bigl(pp_2-p_1^2\bigr)=(a_1-a)pp_1.
\]
\end{theorem}

The square $q^2$ is the local remainder, while $pp_2-p_1^2$ is a discrete
curvature term.  The right side records the change in multiplier, so the
identity supplies a diagnostic for an anti-shadowing argument without
asserting a sign or a global monotonicity.  It is checked as
\lword{Erdos243/DynamicCancellation.lean}{309}{cancellationFree_curvature_square}{the
cancellation-free curvature square}; no global anti-shadowing conclusion is
claimed.

Finally, the reduced denominator transport has a precise realizability gate.
Let $e,e_1$ be centred errors and let $\Delta$ be the Sylvester defect.  If
\[
 hu_1=u-e,
 \qquad he_1=a^2e-\Delta(u-e),
\]
then the transport equality
\[
 h\bigl((a_1-1)u_1+e_1\bigr)
   =a\bigl((a-1)u+e\bigr)
\]
has mismatch which factors through $u-e$.  Consequently, when $u-e\ne0$,
\[
 h\bigl((a_1-1)u_1+e_1\bigr)
   =a\bigl((a-1)u+e\bigr)
 \quad\Longleftrightarrow\quad
 a_1=a^2-a+1+\Delta.
\]
This is the exact
\lword{Erdos243/FeedbackRealizability.lean}{64}{feedback_denominator_transport_iff}{feedback
transport equivalence}, obtained from the factored defect at
\lword{Erdos243/FeedbackRealizability.lean}{15}{feedback_transport_defect}{feedback
transport defect}.  The nonzero-raw-numerator hypothesis is essential: if
$u-e=0$, the factorised mismatch cannot identify $a_1$.

These three identities require orbit-level hypotheses that remain unproved:
eventual absence of cancellation, a fixed sign for the curvature, or cofinal
passage through the feedback gate.  The global mixed-sign and
unbounded-negative regimes of Problem~\#243 remain open.

\begin{theorem}[sequence form]\label{res:eventual}
Let $a,D,C:\N\to\Z$ satisfy $D_{n+1}=\den(a_n,D_n)$ and
$C_{n+1}=\tail(a_n,D_n,C_n)$.  If $E_n=0$ for all sufficiently large $n$ and
$C_{n+1}\ne0$ for all sufficiently large $n$, then $a_{n+1}=a_n^{2}-a_n+1$ for
all sufficiently large $n$.
\end{theorem}

\noindent Formalised as the
\lword{Erdos243/ReciprocalTailRigidity.lean}{1799}{sylvesterNext_eventually_of_centered_zero}{eventual
Sylvester recurrence}.  The proof is routine: take the larger of the two
thresholds and apply Theorem~\ref{res:step} at each later index.  Every later
section is directed at the hypothesis of this theorem, that is, at forcing $E$
to vanish eventually.

The defect identity has a second consequence, which is what makes a single
vanishing error worth having.

\begin{theorem}[zero is absorbing]\label{res:absorb}
Let $a,C,D:\N\to\N$ be an exact natural orbit, so
$C_{n+1}+D_n=a_nC_n$ and $D_{n+1}=a_nD_n$, and let
$E_n=\ctr(a_n,D_n,C_n)$.  Suppose the centring is strict, $|E_n|<C_n$ for
every $n$.  If $E_n=0$ then $E_{n+1}=0$.
\end{theorem}

\begin{proof}
Putting $E_n=0$ in Theorem~\ref{res:defect} gives
$\Delta_n C_{n+1}=-E_{n+1}$, so $C_{n+1}$ divides $E_{n+1}$.  A multiple of
$C_{n+1}$ of absolute value smaller than $C_{n+1}$ is zero, and
$|E_{n+1}|<C_{n+1}$ by hypothesis.
\end{proof}

\noindent Formalised as the
\lword{Erdos243/ReciprocalTailRigidity.lean}{2221}{centeredState_zero_absorbing}{absorption
of a vanishing centred state}, with the companion
\lword{Erdos243/ReciprocalTailRigidity.lean}{2248}{eventually_nonzero_of_zero_absorbing}{contrapositive
along a tail}: if zero is absorbing beyond some index and $E$ does not vanish
eventually, then $E$ is nowhere zero beyond that index.

Absorption changes the shape of the problem.  Either $E$ hits zero once and
stays there, or it is never zero.  Sections~\ref{sec:constant}
and~\ref{sec:periodic} address the special all-negative case only.  The
mixed-sign analysis
in Sections~\ref{sec:bounded} and~\ref{sec:mass} instead separates cofinally
negative errors from an eventually positive tail.

\section{Weighted first crossings of LCM records}\label{sec:lcmrecords}

This section carries the leading result of the note.  It is an exact criterion
for Problem~\ref{res:problem}, stated under the hypotheses of that problem and
no others, and it is reached by a first-crossing argument rather than by a
hypothesis on the negative part.  Only steps setting a new LCM numerator record
contribute, a fixed LCM-scale baseline may be subtracted, and the remaining mass
may be logarithmically discounted.

\paragraph{Evidence.}
The result in this section is an ordinary mathematical proof.  Its finite
arithmetic and charging step are kernel-checked at the pinned revision: the
\lword{Erdos243/LcmRecordExcess.lean}{14}{rise_is_fresh}{freshness of a rise},
the
\lword{Erdos243/LcmRecordExcess.lean}{31}{covered_wall_forces_excess}{covered-wall
excess bound}, the
\lword{Erdos243/LcmRecordExcess.lean}{47}{wall_count_le_excess}{wall count}, and
the
\lword{Erdos243/LcmRecordExcess.lean}{64}{weighted_wall_charge}{weighted wall
charge}.  The global analytic argument, which is the passage from the per-step
charge to divergence of the whole sum, is not formalised.  An exact finite
replay covers $6{,}008$ seed-modulus cases over $6{,}000$ distinct rational
seeds, $26{,}171$ nonterminal transitions and $6{,}008$ terminal states, and
every case reaches centred zero within $15$ steps; finite data of that kind do
not prove universal termination.

\paragraph{What the record coordinates transfer.}
The criterion moves the obligation off the whole orbit.  Only steps setting a
new LCM numerator record contribute, a fixed LCM-scale baseline may be
subtracted, and the remaining mass may be logarithmically discounted, so a
counterexample has to sustain divergence along a sparse subsequence.  The
inherited part of the picture is one half-space.
Proposition~\ref{res:classicalhalfspace} shows that the Erd\H{o}s--Straus
Theorem~3 quantity and Koizumi's Corollary~20(1) quantity carry the sign of
$E_n$ at every index, so each of those published criteria decides the eventual
regime $E_n\ge0$, which Koizumi's Proposition~19(2) closes.  What is added here
is a criterion valid in the complementary regime as well, stated under the
hypotheses of Problem~\ref{res:problem} and no others.  Two exclusions bound
what these coordinates reach.  Proposition~\ref{res:classicalhalfspace}
exhibits cancellation-free blocks on which every classical hypothesis quantity
fails by a fixed number of units, and
Proposition~\ref{res:coprimalitycap} shows that the data used by
Theorem~\ref{res:barrier} and the landing lemmas admit walks with rises of
order $\log u$, so an argument using only those data excludes no negative part
of order $\log C_n$.  A producer that closes the gap has to enter the regime
displayed in Section~\ref{sec:records}, where the negative part is at least
$(1-\delta)\log_2\log_2C_n$ and is $o(C_n)$ at infinitely many indices, and
inside that regime it has to use the congruence $E_n\equiv D_n\pmod{C_n}$
beyond coprimality, or supply a second mechanism that forces landings.

\paragraph{Remaining obligation.}
Derive finiteness of the displayed sum for one admissible weight and one fixed
$B$ from rational denominator feedback, or contradict a nonterminal orbit by
another route.  The logarithmic discount weakens what suffices; it does not
establish it.

Set
\[
 L_n=\operatorname{lcm}(q,a_0,\ldots,a_{n-1}),\quad
 M_n=D_n/L_n,\quad U_n=C_n/M_n,\quad V_n=E_n/M_n.
\]
The rational tail has denominator dividing $L_n$, so $U_n$ and $V_n$ are
integers. With $\rho_n=\gcd(L_n,a_n)$, the exact updates are
\[
 M_{n+1}=M_n\rho_n,\qquad
 \rho_nU_{n+1}=U_n-V_n,\qquad V_n=L_n-(a_n-1)U_n.
\]
Write $R_n=\max_{j\le n}U_j$ and
$\mathcal R=\{n:U_{n+1}>R_n\}$. Centring implies that every sufficiently
late strict rise has $\rho_n=1$: otherwise $U_{n+1}\le3U_n/4$.
Its actual jump is therefore $d_n=U_{n+1}-U_n=-V_n$.

\begin{theorem}[arithmetic weighted-record dichotomy]\label{res:arithmeticrecord}
Let $a_n,L_n,U_n$ be positive integers and $V_n$ integers satisfying
\[
 \begin{gathered}
 L_{n+1}=\operatorname{lcm}(L_n,a_n),\qquad \rho_n=\gcd(L_n,a_n),\\
 \rho_nU_{n+1}=U_n-V_n,\qquad
 V_n=L_n-(a_n-1)U_n,\qquad -U_n\le2V_n.
 \end{gathered}
\]
Let $f:[1,\infty)\to[0,\infty)$ be finite and nonincreasing, with
$\int_1^\infty f(t)\,dt=\infty$. For each fixed integer $B\ge0$,
\[
 \sup_n U_n<\infty
 \quad\Longleftrightarrow\quad
 \sum_{n\in\mathcal R}(-V_n-B)_+f(U_n)<\infty.
\]
\end{theorem}

\begin{proof}
A bounded integer running maximum increases only finitely many times, so
bounded $U_n$ gives a finite sum. Suppose instead that $U_n$ is unbounded.
There are infinitely many record rises, each with $\rho_n=1$. Its digit is
greater than one because $d_n=(a_n-1)U_n-L_n>0$. If $r<s$ are record
indices, then $a_r\mid L_s$ and $\gcd(a_s,L_s)=1$, hence
$\gcd(a_r,a_s)=1$. The record digits are therefore distinct and pairwise
coprime. For any fixed $B\ge1$, choose $B$ of them greater than $B$,
and take $T$ after their indices. These old digits
$m_0,\ldots,m_{B-1}$ all divide $L_T$.

Put $P=\prod_i m_i$ and choose $x$ by the Chinese remainder theorem with
$m_i\mid x+i$. Consider the heights $\tau=x+B+kP$. A first crossing
$U_n\le R_n<\tau\le U_n+d_n$ is a record step. If $d_n\le B$, then
$U_n\in[\tau-B,\tau)$, so some $m_i$ divides $U_n$. It also divides
$L_n$, hence divides $d_n=(a_n-1)U_n-L_n$, contradicting
$0<d_n\le B<m_i$.

If the step first crosses $h\ge1$ such heights, their spacing gives
$(h-1)P<d_n$. With $r=d_n-B\ge1$ and $P\ge B+1$, we have
$d_n=B+r\le Pr$, whence $h\le r$. Monotonicity of $f$ now gives
\[
 \sum_{\substack{\tau\text{ first crossed}\\\text{at step }n}}f(\tau)
 \le(d_n-B)f(U_n).
\]
Each height above $R_T$ has one first crossing. Summing, and comparing each
interval of length $P$ with its left endpoint, yields the finite bound
\begin{equation}\label{eq:weightedcrossing}
 \sum_{\substack{T\le n<N\\n\in\mathcal R}}(-V_n-B)_+f(U_n)
 \ge\frac1P\int_{R_T+P}^{R_N} f(t)\,dt.
\end{equation}
Since $R_n\to\infty$, the right side diverges for every $B\ge1$.
The case $B=0$ follows by domination.
\end{proof}

Only the lower centring bound was used. In particular, neither normalised
vanishing nor a growth estimate for $C_n$ supplies the CRT moduli: the
record digits themselves do so. Normalised vanishing enters in the following
application to the original sequence.

\begin{theorem}[weighted record excess]\label{res:weightedrecord}
Assume the growth and rationality hypotheses of Problem~\ref{res:problem},
and let $f$ be as in Theorem~\ref{res:arithmeticrecord}. The sequence is
eventually Sylvester if and only if, for some integer $B\ge0$,
\[
 \sum_{n\in\mathcal R}(-V_n-B)_+f(U_n)<\infty.
\]
\end{theorem}

\begin{proof}
If the endpoint fails, centred zero is never reached on a sufficiently late
tail. Integrality gives $1/U_n\le|V_n|/U_n=|E_n|/C_n\to0$, so $U_n$ is
unbounded. Theorem~\ref{res:arithmeticrecord}, applied after the centring
threshold, forces divergence for every $B$. Removing a finite prefix changes
only finitely many global record terms. Conversely, a Sylvester tail
telescopes to $x_n=1/(a_n-1)$, so $V_n=0$ eventually.
\end{proof}

For example, $f(t)=1/[t\log(et)]$ gives the sufficient condition
\[
 \sum_{n\in\mathcal R}\frac{(-V_n-B)_+}{U_n\log(eU_n)}<\infty.
\]
Its finite lower bound in~\eqref{eq:weightedcrossing} is
$P^{-1}\log\bigl(\log(eR_N)/\log(e(R_T+P))\bigr)$.
Further fixed iterated logarithmic factors are allowed whenever the integral
still diverges. These are specialisations of one crossing theorem.

The criterion also has an exact expression in the original growth defect.
Put $\gamma_n=a_n^2/a_{n+1}-1$ and $\theta_n=E_n/C_n$. The defect identity
gives
\[
 \gamma_n+\theta_n=
 \frac{(1-\theta_n)(a_n-1+\theta_{n+1})}{a_{n+1}},
 \qquad 0<\gamma_n+\theta_n<3/a_n
\]
eventually. Thus the two nonnegative summands
$U_nf(U_n)(\gamma_n-B/U_n)_+$ and $(-V_n-B)_+f(U_n)$ differ by at most
$3U_nf(U_n)/a_n$. This is summable because
$U_n\le C_n=\exp(o(n))$, $f(U_n)\le f(1)$, and
$a_n\ge\exp(c2^n)$ eventually for some $c>0$.
Consequently Theorem~\ref{res:weightedrecord} is equivalent to finiteness of
\begin{equation}\label{eq:weightedgrowth}
 \sum_{n\in\mathcal R}U_nf(U_n)
 \left(\frac{a_n^2}{a_{n+1}}-1-\frac B{U_n}\right)_+
\end{equation}
for some $B$. The original hypotheses do not currently supply this
finiteness. In particular, termwise convergence to zero is insufficient.
Also, $d_n$ is the actual jump, including any recovery from a drawdown; it
must not be replaced by $R_{n+1}-R_n$ in the crossing proof.


\section{Criteria in the record coordinates}\label{sec:records}

The criterion of Section~\ref{sec:lcmrecords} lives in the LCM coordinates
$U_n=C_n/M_n$, where old divisors persist.  The statements of this section live
in the primitive coordinates, where old divisors can be deleted.  They are
collected here because each is
an equivalence, a dichotomy, or an unconditional exclusion, and because each
sharpens the profile of a counterexample in a direction the earlier sections do
not reach.

\paragraph{Primitive coordinates.}
Put $G_n=\gcd(C_n,D_n)$, $u_n=C_n/G_n$, $v_n=D_n/G_n$, $\tilde e_n=E_n/G_n$
and $h_n=G_{n+1}/G_n$, so that $w_n=a_nu_n-v_n=u_n-\tilde e_n$,
$h_nu_{n+1}=w_n$, $h_nv_{n+1}=a_nv_n$ and $\gcd(u_n,u_{n+1})=1$.  The primitive
error $\tilde e_n$ is signed, and $e_n=-E_n$ of Section~\ref{sec:state} remains
the magnitude in the unreduced coordinates; neither is Koizumi's rational gap
$\varepsilon_n$ of Section~\ref{sec:problem}.  Write $R_n=\max_{k\le n}u_k$ and
$H_n=\max_{j\le n}C_j$ for the two running maxima, $s_n=R_{n+1}-R_n$ for the
true record increment, $d_n=u_{n+1}-u_n$ at a strict rise,
$m_n=(-\tilde e_n)_+$, $\mathcal A_n=R_nm_n/u_n$ for the record-amplified
negative error, $\delta_n=(a_n^2/a_{n+1}-1)_+$ for the growth defect, and
$\ell(x)=\log_2\log_2\max(4,x)$.  The symbol $\mathcal A_n$ is unrelated to the
prefix product $A_n$ of Section~\ref{sec:open}.  A step is \emph{clean} when $h_n=1$.  Call the
standing hypotheses of Problem~\ref{res:problem}, read through Koizumi's
coordinates together with normalised vanishing, the hypotheses~(H), and write
(S) for the conclusion that the sequence is eventually Sylvester.

\paragraph{Source and evidence conventions for this section.}
Each statement below carries its own evidence line.  ``Lean'' means that a
declaration with that statement and a complete proof body is present in the
formal sources.  Several of the modules named below are outside the pinned
public revision~\commitshort; those are marked, and no link is offered for them.
Where a result is an ordinary proof over a Lean core, the Lean part is named
exactly, and the part that is ordinary mathematics is named exactly as well.

\subsection{The record-increment coefficient}

\begin{theorem}[record-excess dichotomy]\label{res:recorddichotomy}
Let $\Theta=\limsup_n(H_{n+1}-H_n)/\ell(H_n)$ on a nonterminating canonical
orbit under~(H).  Then either $G_n$ is unbounded and $\Theta$ is infinite, or
$G_n$ stabilises at a value $g$ and $\Theta\ge g\,v_T/\varphi(v_T)$ for every
late $T$, so that $\Theta>g\ge1$.  Hence $\Theta=0$ or $\Theta>1$, and
$\Theta\le1$ forces~(S).
\end{theorem}

\begin{corollary}[inclusive log-log boundary]\label{res:loglogboundary}
If $\limsup_n(-E_n)_+/\ell(C_n)\le1$ then~(S) holds.  Every counterexample
therefore satisfies $\limsup_n(-E_n)_+/\ell(C_n)>1$ strictly.
\end{corollary}

\paragraph{Mechanism.}
Four steps.  The drawdown ladder $r_n=(-E_n-B_n)_+$ with $B_{n+1}=(B_n+E_n)_+$
and $B_n=H_n-C_n$ gives $0\le r_n\le(-E_n)_+$, so a bound on the record
increment is weaker than a bound on the negative error and is therefore the
better trigger.  A divisor of $C_t$ and $D_t$ at a record locks into every later
record increment, so a pairwise-coprime block of moduli above $B$ dividing $D_T$
contradicts $r_n\le B$ below height $2P+B$.  If $\sup_nG_n$ is infinite then
$\Theta$ is infinite, so a finite $\Theta$ forces a genuinely primitive tail on
which all
multipliers are pairwise coprime.  Pre-sieving the CRT block with $x\equiv0
\pmod{v_N}$ makes a block of length $B$ cost only $\varphi(v_N)v_N^{-1}B$ fresh
moduli, which gives the strict inequality.

\paragraph{Evidence.}
Ordinary proof over the Lean core.  The kernel-checked ingredients are the
\lword{Erdos243/SlowRiseBarrier.lean}{111}{commonDivisor_persists}{persistence of
a common divisor}, the
\lword{Erdos243/SlowRiseBarrier.lean}{51}{exists_consecutiveMultiples_between}{bounded
CRT block}, the
\lword{Erdos243/SlowRiseBarrier.lean}{194}{slowRise_landing}{slow-rise landing
lemma}, and the
\lword{Erdos243/SlowRiseBarrier.lean}{280}{no_slowNegative_of_coprimeBlock}{coprime-block
exclusion}.  The running-maximum persistence clause, the record-increment
trigger, the gcd rescaling and the pre-sieved amplification are ordinary
mathematics.  Finite replay: the identities hold exactly on four orbits, the
$B=20$, $R=6$ CRT reconstruction equals $521026757646$ and excludes all twenty
positions, and the Fermat product identity and the maximal gap ratio $1.8116$ at
$u\ge1000$ against $\sigma^{-1}=2$ hold.

\paragraph{Remaining obligation.}
No upper bound on $\Theta$ is proved.  The branch $\Theta=\infty$ is exactly the
surviving regime.

\subsection{Two exact equivalences under arbitrary cancellation}

\begin{theorem}[record-amplified rigidity]\label{res:recordamplified}
Under~(H), the following are equivalent: (S); $\limsup_n\mathcal A_n<\infty$;
$\limsup_nR_n\delta_n<\infty$.  Each of those two limits superior is $0$ or
$+\infty$.
\end{theorem}

\begin{corollary}[the critical rate]\label{res:criticalrate}
Under~(H) and $\delta_n=O(1/n)$, with no convergence of $n\delta_n$ assumed,
(S) is equivalent to $u_n=O(n)$ and to $(-\tilde e_n)_+=O(1)$.  A
counterexample at the critical rate therefore has $\limsup_nu_n/n=\infty$ and
$\limsup_n(-\tilde e_n)_+=\infty$.
\end{corollary}

\paragraph{Mechanism.}
On a non-Sylvester orbit, arbitrarily late indices carry $B$ distinct primes
above a chosen bound dividing $v_T$, because paid steps have density zero and a
block of consecutive clean steps has pairwise-coprime multipliers all dividing
the final denominator.  A finite alternative between a CRT crossing and a
deletion then forces some $\mathcal A_n$ above any prescribed level.  The passage to the
growth defect uses $|\delta_n-m_n/u_n|\le3/a_n$ and $R_n/a_n\to0$.

\paragraph{Evidence.}
Ordinary proof.  The step that a cancellation at $s$ is still visible at the
first later negative error, in the form $u_t\le u_{s+1}$ and
$u_su_t\le R_t|\tilde e_t|u_{s+1}$, is Lean-checked as
\texttt{recordAmplified\_error\_after\_cancellation} in
\texttt{SaturatedSquareTransport.lean}, which is outside the pinned public
revision~\commitshort.  Exact fixture: $15/134=1/10+1/85+1/5695$, with primitive
steps $(15,134)\to(4,335)\to(1,5695)$ and $\mathcal A_1=15/4$, meeting the
visibility lemma with equality.

\paragraph{Remaining obligation.}
The equivalence relocates the problem into the $\mathcal A_n$ axis.  Nothing
here bounds $\mathcal A_n$ on the canonical orbit.

\begin{theorem}[unit record increments]\label{res:unitrecord}
Under~(H), if $R_{n+1}-R_n\le1$ for all large $n$, then $a_{n+1}=a_n^2-a_n+1$
for all large $n$.  Hence (S) holds if and only if
$\#\{n:R_{n+1}-R_n\ge2\}$ is finite.  No hypothesis is placed on drawdowns, on
record-setting jumps, or on the cancellation factors $h_n$.
\end{theorem}

\paragraph{Mechanism.}
Otherwise $u_n\to\infty$.  A late fresh index $j$ cannot have all its exact prime
powers below $B=R_{j+1}+2$, since $a_j\mid\lcm(1,\ldots,B)$ would give
$\log a_j\le B\log B$ against the doubly exponential growth of $a_j$.  So some
exact prime power $P=p^{k}$ of $a_j$ exceeds $R_s+2$ at $s=j+1$ and divides
$v_s$.  Let $Y$ be the least multiple of $p$ above $R_s$.  At the first $t>s$
with $u_t\ge Y$ the unit increment forces $u_t=Y$ exactly, while the valuation
threshold keeps $P\mid v_t$, so $p\mid u_t$ contradicts primitivity.

\paragraph{Evidence.}
Lean, with an explicit supply hypothesis.  The declarations are
\texttt{unit\_recordIncrement\_cut} and
\texttt{recordIncrementOne\_sylvesterNext\_eventually} in
\texttt{RecordIncrementBarrier.lean}, which is outside the pinned public
revision~\commitshort.  The prime-power supply is a hypothesis of the formal
statement and is discharged at every late index by an ordinary proof, so the
complete advertised endpoint is an ordinary proof with a kernel-checked
principal part, namely the landing contradiction.

\paragraph{Relation.}
This theorem is incomparable with the two-unit record theorem: a record-setting
jump of size at most two bounds $s_n\le2$ rather than $s_n\le1$, and $s_n\le1$
allows arbitrarily large record-setting jumps after a drawdown.  The fixture
$(8,177)\to(7,4071)\to(10,2373393)$ with multipliers $23$ and $583$ has records
$8$ and $10$, increment $2$, jump $3$ and $\gcd(8,10)=2$, which is why the
parity step of the odd cut does not transfer from jumps to increments.

\subsection{Protected epochs and the energy criterion}

\begin{theorem}[protected epoch inequality]\label{res:epochenergy}
Let the orbit satisfy the primitive cocycle with $2|\tilde e_n|<u_n$ from an
index $s$.  Let $p\ge3$ be prime, $Q=p^{\ell}$ divide $v_s$ with $Q\ge16$, put
$L=pQ/2$, and assume $R_s<L/2$.  Let $\tau$ be the first index with $u_\tau\ge
L$, let $J$ be the set of steps in $[s,\tau)$ that first cross at least one odd
multiple of $p$ in $(L/2,L]$, and put $X=\sum_{n\in J}(d_n-2)$.  Then every
$n\in J$ is a clean record step with $d_n\ge3$, and
$pQ\le(8p+8)\lvert J\rvert+4X+8p$.
\end{theorem}

\begin{theorem}[energy criterion]\label{res:energycriterion}
Under~(H), the energy
$\mathcal E=\sum_{n\ \mathrm{record}}\bigl(\mathbf 1_{d_n\ge3}u_n^{-1/2}
+(d_n-2)_+u_n^{-1}\bigr)$ is finite if and only if~(S) holds; and
$\sum_{n\ \mathrm{record}}(d_n-2)_+u_n^{-1/2}$ is finite if and only if~(S)
holds.
\end{theorem}

\paragraph{Mechanism.}
Below $L$ the numerator satisfies $w_n<\tfrac32L<p^{\ell+1}$, so the prime power
persists and $p\nmid u_t$ throughout $[s,\tau]$.  Each barrier above $R_s$ is
first crossed at a record step, where a jump of at most two would require either
the forbidden landing or a pair of even endpoints.  Barriers first crossed at one
step lie in an interval of length $d_n$ with spacing $2p$, which gives the
per-step capacity $1+d_n/(2p)$, and the barrier count is at least $Q/8-1$.  For
the criterion, an odd prime power of the required size is available at every late
index, so each late window carries energy at least $1/16$, and the tail of a
convergent series tends to zero.

\paragraph{Evidence.}
The integer epoch inequality is Lean, as
\texttt{protected\_epoch\_energy\_integer} with its three counting lemmas in
\texttt{ProtectedEpochEnergy.lean}, which is outside the pinned public
revision~\commitshort.  The real-valued corollary, the odd prime-power supply at
every late index, and the global criterion are ordinary proofs.  This discharges
a quantitative record-excess target that an earlier wave had left open, and it
uses one protected prime power rather than a product of moduli.

\subsection{A slow negative part, and the original coordinates}

\begin{theorem}[slow negative part]\label{res:slownegative}
Let $(a,C,D)$ be an exact natural orbit with $a_n>1$, $C_n>0$, $D_0\ge1$, under
normalised vanishing.  Suppose that for some $\delta\in(0,1)$ and all large $n$
with $E_n<0$ one has $-E_n\le(1-\delta)\ell(C_n)$.  Then $E_n=0$ for all large
$n$, and $a_{n+1}=a_n^2-a_n+1$ for all large $n$.
\end{theorem}

\noindent Theorem~\ref{res:bounded} is the case of a constant bound, since a
constant is eventually below $(1-\delta)\ell(C_n)$.

\paragraph{Mechanism.}
Normalised vanishing gives $C_{n+1}<\tfrac32C_n$, hence
$\ell(C_n)\le\log_2n+O(1)$, and $a_n$ grows at least geometrically.  The
multiplier overlap $\gcd(a_n,D_n)$ is bounded by the next negative magnitude and
therefore by $\ell(C_n)+1$.  Splitting $a_n$ into the part coprime to $D_n$ and
the burnt part, a poor index supplies a burn prime with valuation above
$\sqrt n$, and such primes are injective across a window, so a window
$[T,2^T]$ has at most $T+O(1)$ poor indices.  The first $B$ rich indices give
$B$ pairwise-coprime moduli above any prescribed size, and the block form of the
landing lemma applies.

\paragraph{Evidence.}
Ordinary proof over the Lean core.  The kernel-checked ingredients are the
\lword{Erdos243/SlowRiseBarrier.lean}{111}{commonDivisor_persists}{persistence of
a common divisor}, the
\lword{Erdos243/SlowRiseBarrier.lean}{162}{multiplierOverlap_dvd_laterCenteredState}{overlap
transport}, the
\lword{Erdos243/SlowRiseBarrier.lean}{194}{slowRise_landing}{landing lemma} and
the
\lword{Erdos243/SlowRiseBarrier.lean}{280}{no_slowNegative_of_coprimeBlock}{coprime-block
exclusion}.  The burn count and the choice of constants are ordinary
mathematics.  Corollary~\ref{res:loglogboundary} supersedes this theorem on two
axes, the drawdown-corrected observable and the inclusive constant, and this
theorem is retained because its hypothesis is stated directly on $E_n$.

\begin{theorem}[the Erd\H{o}s--Straus quantity, bounded form]\label{res:strausbounded}
Let $a_{n+1}/a_n^2\to1$ and $\sum1/a_n\in\Q$, and put
\[
 Q_n=\frac{a_1a_2\cdots a_{n-1}}{a_n}
     \left(\frac{a_n^2}{a_{n+1}}-1\right).
\]
If $\limsup_nQ_n<\infty$ then the sequence is eventually Sylvester.  If for
some $\delta>0$ and all large $n$ one has
$Q_n\le(1-\delta)\lambda\,\ell(a_1\cdots a_{n-1}/a_n)$, where $\lambda>0$
depends only on the sequence, the same conclusion holds.
\end{theorem}

\paragraph{Attribution.}
Erd\H{o}s and Straus require $\limsup_nQ_n\le0$ in the corresponding criterion,
with the least common multiple in place of the product~\cite[Theorem~3,
p.~132]{erdosstraus1964}; Koizumi's Corollary~20(1) uses the product
form~\cite[Cor.~20(1), p.~12]{koizumi2025}.  The statement above is the same
Lean-checked endpoint as Theorem~\ref{res:bounded}, transported into that
language, and its content over the published criteria is the passage from
$\le0$ to $<\infty$.  It is the only place in this note where the machine-checked
theorem is written in the vocabulary of the 1964 paper.

\paragraph{Evidence.}
Ordinary proof composing Theorem~\ref{res:bounded} and
Theorem~\ref{res:slownegative} with Koizumi's bridge.  The identity used is
$Q_n=-(B_{n-1}/c_1)\varepsilon_n+(B_{n-1}/c_1)\beta_nc_{n+1}$ with
$\beta_k=O(1/a_k)$, whose second term tends to zero because $c_{n+1}\le
c_1(3/2)^n$ while $a_n$ is doubly exponential.

\begin{theorem}[the $1/n$ threshold]\label{res:onethreshold}
Let $a_{n+1}/a_n^2\to1$ and $\sum1/a_n\in\Q$.  If
$\limsup_n n\,(a_n^2/a_{n+1}-1)_+<1$, then $a_{n+1}=a_n^2-a_n+1$ for all large
$n$.
\end{theorem}

\noindent The constant $1$ is the boundary of this method.  The template
$\tilde e_n\equiv-1$ with $u_n=c+n$ has $n(-\tilde e_n)/u_n\to1$, meets
every hypothesis of the Erd\H{o}s--Straus criterion except the sign, and is
excluded only by Theorem~\ref{res:constant}.  This sharpens the sufficient rate
$a_n^2/a_{n+1}=1+o(1/n)$ recorded by Koizumi~\cite[Remark~21,
p.~13]{koizumi2025} to an explicit constant.  Ordinary proof; the endgame is
Koizumi's Proposition~19(2).

\subsection{What the method cannot reach}

\begin{proposition}[the classical criteria decide one half-space]\label{res:classicalhalfspace}
The overlap debt divides the tail gcd, so $\kappa'_n=G_n/M_n$ is a positive
integer and $E_n/M_n=\kappa'_n\tilde e_n$ is an integer carrying the sign of
$E_n$.  Consequently the Erd\H{o}s--Straus
Theorem~3 quantity and Koizumi's Corollary~20(1) quantity carry the sign of
$E_n$ at every index, and each of those criteria decides exactly the half-space
$E_n\ge0$ eventually, which Koizumi's Proposition~19(2) closes.  Moreover
$\limsup_nQ_n/M_n<\infty$ holds if and only if $(-E_n)_+=O(M_n)$; that
condition is equivalent to $\limsup_nQ_n<\infty$ on a cancellation-free tail,
and every record of $u$ is set by a clean step, at which $M$ does not grow.
\end{proposition}

\paragraph{Evidence.}
Ordinary proof over exact identities, with an exact finite replay on $4{,}472$
frames and $610$ clean triples, zero failures.  The fixture family with
$C_n=c+nm$, $D_n=(a_n-1)C_n-m$ and $a_n=1+C_{n+1}k_n$ realises, at
$(c,m,h)=(10,3,6)$, $(13,3,6)$ and $(100003,3,5)$, arbitrarily long blocks on
which every classical hypothesis quantity equals $-E_n=3$ at every index while
$G_n=\kappa'_n=M_n=1$: the classical hypotheses fail by exactly $m$ units on a
cancellation-free block.

\paragraph{Reading.}
This is a negative result about the classical toolkit and about the note's own
LCM-weighted question.  It proves nothing about the surviving regime.  A third
classical quantity, from Tijdeman and Yuan's 2002 Theorem~4.1, enters the same
computation; it is not cited here because this note carries no verified
bibliographic record of that paper.

\begin{proposition}[cap on the coprimality route]\label{res:coprimalitycap}
Let $m_0<m_1<\cdots$ be pairwise coprime integers at least $2$ with
$\sum_i1/m_i\le1/2$.  For every $x\ge1$ and every $L>2^{k+1}$, where
$k=\#\{i:m_i\le x+L\}$, the window $[x,x+L)$ contains an integer coprime to
every $m_i$.  Consequently there is a sequence $u_n\to\infty$ coprime to every
$m_i$ with $u_{n+1}-u_n\le2^{k(u_n)+1}+1$, which for $m_i\ge2^{2^i}$ is a rise
of order $\log u_n$.  For pairwise-coprime $m_j$ with $\ell(m_j)=j+O(1)$ and
$\sigma=\prod_j(1-1/m_j)>0$, the whole-modulus-avoiding walk has
$\limsup_n(u_{n+1}-u_n)/\ell(u_n)=\sigma^{-1}$ exactly, and the Fermat family
$m_j=2^{2^j}+1$ gives the constant exactly $2$.
\end{proposition}

\paragraph{Reading.}
Theorem~\ref{res:barrier} and the landing lemmas use about the multipliers only
that they are pairwise coprime, that they divide later denominator states, and
that a modulus dividing the current numerator is impossible at a landing.
Proposition~\ref{res:coprimalitycap} shows those data admit walks with rises of
order $\log u$, while the CRT block forces rises only of order $\ell(u)$.  No
argument using only those data can exclude negative parts of order $\log C_n$,
and the constant $1$ in Corollary~\ref{res:loglogboundary} is the CRT constant
rather than a feature of the problem.  Ordinary proof, with a numerical witness
walk whose maximal rise against Sylvester's numbers as moduli is $12$ up to
height $2\cdot10^6$.

\paragraph{The exact remaining obstacle.}
Every criterion of this section leaves the same regime, which is
\[
 (1-\delta)\ell(C_n)\ \le\ -E_n\ =\ o(C_n)
 \qquad\text{at infinitely many }n,
\]
with divergent normalised negative mass, unbounded record-amplified error, and
infinitely many clean primitive record jumps of size at least three.  Any
argument entering it must use the congruence $E_n\equiv D_n\pmod{C_n}$ beyond
coprimality, or supply a second mechanism that forces landings.

\section{Descent when the error is nonnegative}\label{sec:descent}

The hypothesis of the theorem below is the update law of
Proposition~\ref{res:update} read inside $\N$: the error is nonnegative at
every index, so the tail state never rises.  This is the easy case, and it is
settled in three lines.

\begin{theorem}[descent]\label{res:descent}
Let $C,E:\N\to\N$ satisfy $C_{n+1}+E_n=C_n$ for every $n$.  Then $E_n=0$ for
all sufficiently large $n$.
\end{theorem}

\begin{proof}
A one-line descent.  The relation forces $C_{n+1}\le C_n$, so $C$ is a nonincreasing sequence of
natural numbers and is eventually equal to its minimum; beyond that index
$C_{n+1}=C_n$, and the relation gives $E_n=0$.
\end{proof}

\noindent Formalised as the
\lword{Erdos243/ReciprocalTailRigidity.lean}{1837}{centeredState_eventually_zero}{stabilisation
of the nonnegative error}, on the
\lword{Erdos243/ReciprocalTailRigidity.lean}{1819}{antitone_nat_eventually_constant}{eventual
constancy of a nonincreasing sequence}.

The hypothesis is exactly one-sidedness: $C$ and $E$ take values in $\N$.  If
$E_n<0$ at infinitely many indices then $C$ rises there, nothing descends, and
the argument is unavailable.  That is the whole difficulty, and the rest of
this note is directed at it.

\section{Excluding a constant negative magnitude}\label{sec:constant}

Suppose the error is negative with a constant magnitude, $E_n=-m$ for a fixed
$m>0$ and every $n$.  By
Proposition~\ref{res:update} the state then climbs by exactly $m$ each step,
so $C_n=c+nm$ with $c=C_0$, and the definition of the centred state becomes a
\emph{shape equation}
\[
 D_n+m=(a_n-1)\,(c+nm) .
\tag{5.1}\label{eq:shape}
\]
Together with $D_{n+1}=a_nD_n$ this is a closed system in $(a,D)$, and it has
no solutions.  The following instance shows how the failure happens; the
theorem then shows that it cannot be avoided.

\begin{example}[the shape equation running until it fails]\label{ex:shape}
Take $m=c=1$, so that $C_n=1+n$ and the shape equation~\eqref{eq:shape} reads
$D_n+1=(a_n-1)(1+n)$, and take $D_0=1$.  Each step is now determined: at $n=0$
the equation reads $2=(a_0-1)\cdot1$, so $a_0=3$, and $D_1=a_0D_0=3$; at $n=1$
it reads $4=(a_1-1)\cdot2$, so $a_1=3$, and $D_2=a_1D_1=9$; at $n=2$ it reads
$10=(a_2-1)\cdot3$, which has no integer solution, and the orbit stops.
Longer prefixes occur for other starting values: the finite search recorded in
Appendix~\ref{app:residue} found none longer than $17$ among the initial
states below $5000$.
\end{example}

\begin{theorem}[no constant negative magnitude]\label{res:constant}
There is no pair of sequences $a,D:\N\to\N$ with $a_n\ge2$ for all $n$
satisfying $D_{n+1}=a_nD_n$ and~\eqref{eq:shape}, for any $m>0$ and any
$c$.  The same holds if the shape only begins at some index.
\end{theorem}

\begin{proof}[Proof when $c$ and $m$ are coprime]
Assume first $\gcd(c,m)=1$.  Two facts about the multipliers do all the work,
and the difficulty lies in the fact that they concern the same finite set of
primes from opposite directions: the first makes every multiplier meet that
set, the second lets each prime of the set meet at most one multiplier.

\emph{Every $a_j$ shares a prime with $m$.}  Suppose $\gcd(a_j,m)=1$.  Then
$m$ is invertible modulo $a_j$, so some $n$ has $c+nm\equiv0$, that is
$a_j\mid C_n$; and $a_j\mid D_n$ for every $n>j$, since $D$ is multiplicative
with $a_j$ among its factors.  Choosing such an $n$ beyond $j$
and reading~\eqref{eq:shape} modulo $a_j$ gives $a_j\mid m$, contradicting
$\gcd(a_j,m)=1$ and $a_j\ge2$.

\emph{A prime divisor of $m$ occurs in at most one $a_j$.}  Suppose $p\mid m$
and $p\mid a_j$.  For $n>j$ we have $p\mid D_n$, so~\eqref{eq:shape} gives
$(a_n-1)C_n\equiv0 \pmod p$.  Now $C_n=c+nm\equiv c$, and $p\nmid c$ because
$p\mid m$ and $\gcd(c,m)=1$; hence $p\mid a_n-1$, so $p\nmid a_n$.  Thus $p$
cannot divide any later multiplier.

The two facts are incompatible.  Each of the infinitely many multipliers needs
a prime divisor of $m$, and each prime divisor of $m$ serves at most one
multiplier, while $m$ has only finitely many prime divisors.
\end{proof}

\noindent The case $\gcd(c,m)=1$ is the
\lword{Erdos243/ReciprocalTailRigidity.lean}{147}{no_scalePrimitiveConstantNegative_orbit}{scale-primitive
exclusion}.  The special case $m=c=1$, worked through in
Example~\ref{ex:shape}, where $m$ has no prime divisors at all and the first
fact is immediately contradictory, is recorded separately as the
\lword{Erdos243/ReciprocalTailRigidity.lean}{100}{no_normalizedConstantNegative_orbit}{normalised
exclusion}.

\begin{proof}[Removing the scale]
For general $c$, put $g=\gcd(c,m)$.  Equation~\eqref{eq:shape} shows
$g\mid D_n$ for every $n$, so dividing $D$, $c$ and $m$ by $g$ leaves a system
of the same shape with coprime data, which the previous case excludes.
\end{proof}

\noindent Formalised as the
\lword{Erdos243/ReciprocalTailRigidity.lean}{286}{no_constantNegative_orbit}{exclusion
at every scale}, and the eventual form, obtained by shifting the orbit to the
first constant index, as the
\lword{Erdos243/ReciprocalTailRigidity.lean}{350}{no_eventuallyConstantNegative_orbit}{eventual
exclusion}.

Theorem~\ref{res:constant} excludes a centred state that is eventually constant
and negative, at every magnitude and every scale.  It says nothing about a
negative state whose magnitude changes, and the argument genuinely uses
constancy: the finiteness of the set of prime divisors of $m$ is what the
pigeonhole uses, and a varying magnitude presents a fresh set at each index.

\section{Excluding a periodic negative magnitude}\label{sec:periodic}

The next case allows the magnitude to vary, provided it repeats.  Write
$e_n=-E_n>0$ for the magnitude, as in Section~\ref{sec:state}, so that
$C_{n+1}=C_n+e_n$; suppose $e$ has period $h$, meaning $e_{n+h}=e_n$ for every
$n$, and that the state gains a fixed \emph{drift} $M>0$ over one period,
meaning $C_{n+h}=C_n+M$ for every $n$.

At $h=1$ this says that the magnitude is constant and that $M$ is its value,
which is the situation of Section~\ref{sec:constant};
Theorem~\ref{res:constant} already excludes it, and does so without the extra
hypothesis $e_n<a_n$ imposed below.  The new content is $h\ge2$.

Section~\ref{sec:constant} could run its pigeonhole on the prime divisors of
the single number $m$.  Here the magnitude changes inside a period, and the
number available to run it on is the drift $M$.  Three lemmas replace the two
facts of Section~\ref{sec:constant}.  The first is
that multipliers persist: every $a_j$ divides every strictly later denominator
state, the
\lword{Erdos243/ReciprocalTailRigidity.lean}{382}{base_dvd_denState_later}{persistence
of a multiplier}.  The second is a \emph{lock}: a prime already dividing $D$ and also
dividing the current multiplier is forced into the next tail state, the
\lword{Erdos243/ReciprocalTailRigidity.lean}{402}{divisorLocksTailSucc}{one-step
lock}, and once present in both $D$ and $C$ it stays in every later $D$, $C$
and magnitude, the
\lword{Erdos243/ReciprocalTailRigidity.lean}{429}{divisorLock_persists}{persistence
of a lock}.  The third transports a lock backwards through the period: if a
divisor of the drift $M$ occurs in two multipliers, it is a common divisor of
$C_0$ and of every magnitude, the
\lword{Erdos243/ReciprocalTailRigidity.lean}{494}{repeatedDivisor_forces_periodicCommonScale}{repeated-divisor
transport}, using that both the period relation and the drift iterate over any
number of periods, the
\lword{Erdos243/ReciprocalTailRigidity.lean}{461}{periodic_add_mul}{iterated
period} and the
\lword{Erdos243/ReciprocalTailRigidity.lean}{476}{phaseDrift_add_mul}{iterated
drift}.

\begin{theorem}[no periodic negative magnitude]\label{res:periodic}
Let $a,D,C,e:\N\to\N$ with $a_n\ge2$, $e_n>0$ and $e_n<a_n$ for every $n$,
satisfying
\[
 D_{n+1}=a_nD_n,\qquad C_{n+1}=C_n+e_n,\qquad D_n+e_n=(a_n-1)C_n ,
\]
and suppose $e_{n+h}=e_n$ and $C_{n+h}=C_n+M$ for some $h>0$ and $M>0$.  This
is impossible.
\end{theorem}

\begin{proof}
Strong induction on the drift $M$.  The difficulty lies in the fact that a
repeated prime divisor of $M$ need not contradict anything directly.  It
instead forces a common scale on the whole orbit, which can be divided out, and
the induction runs on the drift that the division reduces.

Suppose first that no prime divisor of $M$ is a common divisor of $C_0$ and of
every magnitude.  Repeated-divisor transport then applies in the
contrapositive: a prime divisor of $M$ occurring in two of the multipliers
would be exactly such a common divisor, so each prime divisor of $M$ occurs in
at most one multiplier.  Against this, the argument of
Section~\ref{sec:constant} run phase by phase (along a fixed residue class
modulo $h$ the magnitude is constant and $C$ advances by $M$ at each period, so
the multipliers of that phase meet an arithmetic progression exactly as in
Section~\ref{sec:constant}) makes every multiplier share a prime with $M$.
Since $M$ has only finitely many prime divisors and there are infinitely many
multipliers, pigeonhole gives a contradiction.

Otherwise some prime $p$ divides $M$, divides $C_0$, and divides every
magnitude.  Then $p\mid C_n$ for every $n$, since $C_{n+1}=C_n+e_n$ and both
summands on the right are divisible by $p$, and the shape equation
$D_n+e_n=(a_n-1)C_n$ then gives $p\mid D_n$ for every $n$.  Divide $D$, $C$ and
$e$ by $p$.  The multipliers are untouched, so $a_n\ge2$ and $e_n<a_n$ persist
and the magnitudes stay positive; the three recurrences are homogeneous in
$(D,C,e)$ and so survive; the period is still $h$; and the drift becomes
$M/p<M$.  The inductive hypothesis applies.
\end{proof}

\noindent The first of the two cases above is the
\lword{Erdos243/ReciprocalTailRigidity.lean}{548}{no_phasePrimitivePeriodicNegative_orbit}{phase-primitive
exclusion}, the induction is the
\lword{Erdos243/ReciprocalTailRigidity.lean}{690}{no_periodicNegative_orbit}{exclusion
at every scale}, and the eventual form is the
\lword{Erdos243/ReciprocalTailRigidity.lean}{796}{no_eventuallyPeriodicNegative_orbit}{eventual
exclusion}.

The bound $e_n<a_n$ is a genuine restriction and not a normalisation.  It is
the range in which a multiplier cannot divide its own phase magnitude, which
is what the lock argument needs.  It holds in the intended reading, where
$e_n$ is a centred representative and $a_n$ grows quadratically, but it is
assumed here and not derived.  The drift hypothesis $M>0$ is also used: a
periodic magnitude with zero drift would be identically zero, which is the case
already settled by Section~\ref{sec:descent}.

\begin{remark}
On the canonical orbit of~\cite{koizumi2025}, read through the dictionary of
Section~\ref{sec:problem}, the bound $e_n<a_n$ holds eventually.  Lemma~15(2)
there gives the centring range $-C_n/2\le E_n<C_n/2$, so $e_n=|E_n|\le C_n/2$ at
a negative index, and the proof of Lemma~15(3) gives $C_{n+1}\le\tfrac32C_n$
~\cite[p.~9]{koizumi2025},
hence the uniform bound $C_n\le C_0(3/2)^{n}$.  Against this, summability of
$\sum1/a_n$ forces
$a_n\to\infty$, and the growth hypothesis $a_n^{2}/a_{n+1}\to1$ gives
$a_{n+1}\ge a_n^{2}/2$ for all large $n$, hence $a_{n+1}\ge2a_n$ once
$a_n\ge4$; so the multipliers grow at least geometrically with ratio $2$ and
outgrow the bound on $C_n$, making $e_n/a_n$ eventually small.  Two caveats.
Summability is used for this step, so the estimate is not available from the
state recurrence alone; and the conclusion is eventual and not universal,
whereas Theorem~\ref{res:periodic} imposes $e_n<a_n$ at every index, so it is
the eventual form of the exclusion cited above that applies.  Under the
periodicity hypothesis of Theorem~\ref{res:periodic} the estimate is in any
case unnecessary, since a periodic magnitude is bounded while the multipliers
tend to infinity.
\end{remark}

\section{Bounded rise cannot remain coprime to fresh moduli}\label{sec:barrier}

Periodicity is still a strong assumption.  Removing it needs a different kind
of obstruction, and the one used here is a counting argument about where an
integer sequence with bounded upward steps must land.

\begin{lemma}[shifted blocks of consecutive multiples]\label{res:crt}
Let $m_0,\ldots,m_{B-1}$ be pairwise coprime and at least $2$.  For every bound
there is a $t$ beyond it with $m_i\mid t+i$ for each $i<B$.
\end{lemma}

\begin{proof}
Standard Chinese remaindering: solve $t\equiv-i \pmod{m_i}$ simultaneously,
then add multiples of $\prod_i m_i$ to pass the bound.
\end{proof}

\noindent Formalised as the
\lword{Erdos243/ReciprocalTailRigidity.lean}{833}{exists_shifted_consecutiveMultiples}{shifted
consecutive multiples}.

\begin{example}[a block of three consecutive multiples]\label{ex:crt}
Take $B=3$ and $(m_0,m_1,m_2)=(3,4,5)$.  The congruences $t\equiv0\pmod3$,
$t\equiv3\pmod4$ and $t\equiv3\pmod5$ hold exactly when $t\equiv3\pmod{60}$.
At $t=3$ the block is $3,4,5$ with $3\mid3$, $4\mid4$ and $5\mid5$; at $t=63$
it is $63,64,65$ with $3\mid63$, $4\mid64$ and $5\mid65$.  A sequence of
natural numbers that starts at $0$, tends to infinity and rises by at most $3$
at each step cannot step over the block $\{3,4,5\}$: it has a first value at
least $3$, that value is at most $5$, and every member of $\{3,4,5\}$ is
divisible by one of the three moduli.
\end{example}

\begin{theorem}[bounded rise cannot remain coprime to fresh moduli]\label{res:barrier}
Let $u:\N\to\N$ tend to infinity with $u_{n+1}\le u_n+B$ for a fixed $B>0$.
Then $u$ cannot remain coprime to infinitely many fresh pairwise coprime
moduli: there is no family of pairwise coprime $m_i\ge2$, one for each index,
such that $\gcd(m_i,u_t)=1$ whenever $i<t$.
\end{theorem}

\begin{proof}
Take $B$ of the moduli, from indices beyond any chosen point, and use
Lemma~\ref{res:crt} to find $t$ far out with the $i$th of them dividing $t+i$.
The interval $[t,t+B)$ has length $B$, and $u$ climbs by at most $B$ per step
while tending to infinity, so some $u_s$ lands in it.  Then $u_s=t+i$ for some
$i<B$, and the $i$th modulus divides $u_s$.  Since that modulus is at least
$2$, this contradicts $\gcd(m_i,u_s)=1$.
\end{proof}

\noindent Formalised as the
\lword{Erdos243/ReciprocalTailRigidity.lean}{897}{no_boundedRise_of_tailAvoidance}{bounded-rise
barrier}.  The key point is that the two hypotheses pull against each other
exactly: divergence forces $u$ to travel arbitrarily far, and bounded rise
forbids it from stepping over a window of width $B$.  Uniformity of the bound
is what the proof uses, and it is the hypothesis that cannot be relaxed by this
argument: a rise bounded only along a subsequence leaves the intervening steps
free to jump the window.  We have not looked for a sequence of unbounded rise
remaining coprime to every fresh modulus.  The statement is elementary and we are not aware of a
prior published form of it, but we would not be surprised if one exists.

The moduli are indexed by the same set as the sequence, and the $i$th of them
is asked to divide no value $u_t$ with $t>i$; nothing at all is asked of it at
or before its own index.  That asymmetry is what makes the hypothesis available
in the application below, where $m_i$ is the $i$th multiplier and only the
later numerators are known to be coprime to it.

The barrier applies to the problem because a \emph{reduced} exact tail already
carries pairwise coprime moduli.  Call $(u,v)$ a reduced exact tail with
multipliers $a$ when $\gcd(u_n,v_n)=1$,
\[
 u_{n+1}+v_n=a_nu_n,\qquad v_{n+1}=a_nv_n .
\]
Here $u$ is the numerator and $v$ the denominator: the two recurrences are
those of Section~\ref{sec:state}, with $u$ in the role of the tail state $C$
and $v$ in that of the denominator state $D$, and with coprimality imposed at
every index.

\begin{proposition}[reduced tails are pairwise coprime]\label{res:reduced}
In a reduced exact tail, $a_n$ is coprime to $v_n$; the multipliers at distinct
indices are pairwise coprime; and every earlier multiplier is coprime to every
later numerator.
\end{proposition}

\noindent These are the
\lword{Erdos243/ReciprocalTailRigidity.lean}{983}{reducedStep_coprime_currentFactor}{step
coprimality}, the
\lword{Erdos243/ReciprocalTailRigidity.lean}{1010}{reducedTail_pairwiseCoprime}{pairwise
coprimality}, and the
\lword{Erdos243/ReciprocalTailRigidity.lean}{1024}{reducedTail_wholeModulusAvoidance}{whole-modulus
avoidance}.  Feeding them to Theorem~\ref{res:barrier} gives the form used
later: there is no reduced exact tail whose numerator tends to infinity with a
uniformly bounded upward increment, the
\lword{Erdos243/ReciprocalTailRigidity.lean}{1035}{no_boundedRise_reducedTail}{reduced
bounded-rise exclusion}, together with its eventual form, the
\lword{Erdos243/ReciprocalTailRigidity.lean}{1060}{no_eventuallyBoundedRise_reducedTail}{eventual
reduced exclusion}.

\begin{example}[Sylvester's sequence as a reduced exact tail]\label{ex:reduced}
Put $u_n=1$ and $v_n=D_n=1,2,6,42,1806,\ldots$ as in
Example~\ref{ex:sylvester}, with multipliers $a_n=v_n+1=2,3,7,43,1807,\ldots$.
Then $\gcd(u_n,v_n)=1$, $u_{n+1}+v_n=1+v_n=a_nu_n$ and $v_{n+1}=a_nv_n$, so
this is a reduced exact tail, and Proposition~\ref{res:reduced} returns the
classical fact that Sylvester's numbers are pairwise coprime.  Its numerator is
constant, so it satisfies every hypothesis of the exclusion just stated except
divergence; since the tail exists, that hypothesis cannot be dropped.
\end{example}

An arbitrary orbit of Section~\ref{sec:state} need not be reduced, so one more step is needed
before the barrier can be applied to it: the common factor must stop changing.
We call a set of indices \emph{cofinal} when it is infinite, and say that a
property holds \emph{cofinally} when the set of indices at which it holds is
cofinal.  The step below is the only one connecting the orbit of
Section~\ref{sec:state} to the reduced tails of
Proposition~\ref{res:reduced}, and it is where the cofinal boundedness in the
proof of Theorem~\ref{res:bounded} is used.

\begin{proposition}[the tail gcd stabilises]\label{res:gcdstab}
Let $(a,D,C)$ be an exact natural orbit, that is, a triple of $\N$-valued
sequences with $C_{n+1}+D_n=a_nC_n$ and $D_{n+1}=a_nD_n$, whose centred state
$E_n=\ctr(a_n,D_n,C_n)$ is negative cofinally, with magnitudes bounded along
that cofinal set.  Then
$\gcd(C_n,D_n)$ is eventually constant, and beyond that index the orbit divided
by the stable gcd is a reduced exact tail.
\end{proposition}

\begin{proof}
The tail gcd $\gcd(C_n,D_n)$ divides its successor, the
\lword{Erdos243/ReciprocalTailRigidity.lean}{1424}{tailGcd_dvd_succ}{gcd
monotonicity}, so the tail gcds form a positive divisibility chain.  At an
index where the centred state is negative the tail gcd is exactly
$\gcd(C_n,e_n)$, the
\lword{Erdos243/ReciprocalTailRigidity.lean}{1581}{tailGcd_eq_gcd_negativeMagnitude}{gcd
at a negative index}, and is therefore at most the magnitude there; so the
chain is bounded along the cofinal set of negative indices.  A positive
divisibility chain whose values are bounded along a cofinal set is eventually
constant, the
\lword{Erdos243/ReciprocalTailRigidity.lean}{1599}{dvdChain_eventuallyConstant_of_cofinally_bounded}{chain
stabilisation}.  Hence cofinally bounded negative magnitudes make the tail gcd
stabilise, the
\lword{Erdos243/ReciprocalTailRigidity.lean}{1630}{tailGcd_eventuallyConstant_of_cofinally_boundedNegative}{gcd
stabilisation}, and beyond that point dividing by the stable gcd leaves the
numerator and denominator coprime, that is, a reduced exact tail.
\end{proof}

\noindent The stabilisation of the tail gcd is the checked statement cited at
the end of the proof; the last clause, that the divided orbit is a reduced
exact tail, is exposition and carries no separate label.  Boundedness is
required only \emph{along a cofinal set}.  That is what the situation gives:
the tail gcd is
controlled only at the negative indices, and cofinally many of them is all the
divisibility chain needs.

Normalised vanishing by itself gives a different, strictly weaker conclusion.
For an exact natural orbit with $C_n>0$, put
\[
 G_n=\gcd(C_n,D_n),\qquad
 \Gamma(N)=\#\{0\le j<N:G_j<G_{j+1}\}.
\]

\begin{proposition}[normalised vanishing makes strict gcd changes sparse]
\label{res:gcdsparse}
Assume
\[
  (\forall K\ge1)(\exists N)(\forall n\ge N)\qquad
  K|E_n|<C_n.
\]
Then, for every $K\ge1$, eventually $K\Gamma(N)<N$.  Moreover, for every
starting bound $B$ and block length $L$, some $n\ge B$ satisfies
\[
  G_n=G_{n+1}=\cdots=G_{n+L}.
\]
\end{proposition}

\noindent The first assertion is the
\lword{Erdos243/ReciprocalTailRigidity.lean}{2180}{tailGcd_strictGrowthCount_sublinear_of_normalizedVanishes}{sublinear
strict-growth theorem}; the second is the
\lword{Erdos243/ReciprocalTailRigidity.lean}{2196}{tailGcd_exists_arbitrarilyLate_constBlock_of_normalizedVanishes}{arbitrarily
late constant-block theorem}.  The proof first converts normalised vanishing
into subexponential growth of $C_n$, then charges every strict divisibility
increase of $G_n$ against a finite power-of-two budget.  Sparse strict changes
force long finite gaps between them.

The quantifiers matter.  Proposition~\ref{res:gcdstab} uses cofinally bounded
negative magnitudes and yields eventual constancy.  Proposition~\ref{res:gcdsparse}
uses only normalised vanishing and yields arbitrarily late constant blocks of
each prescribed finite length.  It does not yield one infinite constant tail,
and it does not exclude cofinally unbounded negative excursions.

\section{Excluding a bounded negative part}\label{sec:bounded}

Assembling the previous section gives Theorem~\ref{res:bounded}, the endpoint
of the barrier argument.  Theorem~\ref{res:mass} in the next section uses a
different finiteness hypothesis.  The two results are incomparable:
Theorem~\ref{res:mass} assumes only the state recurrence with $C_n>0$, while
hypothesis~(5) below does not imply
finite \emph{normalised negative mass}, by which we mean convergence of
$\sum_n(-E_n)_+/C_n$, where $x_+=\max(x,0)$, since a constant negative part
$-b$ gives $C_n\sim bn$ and hence a divergent sum.  The hypotheses of
Theorem~\ref{res:bounded} are listed
in full, because two of them are analytic inputs
that are assumed here rather than derived from the growth condition of
Problem~\ref{res:problem}.

Informally, hypotheses~(5) and~(6) below pull against each other.  By
Proposition~\ref{res:update}, hypothesis~(5) caps the rise of the tail state at
$B$ per step, while hypothesis~(6), once the error is nowhere zero, forces the
tail state to diverge.  Proposition~\ref{res:gcdstab} makes the orbit reduced,
and Theorem~\ref{res:barrier} then converts the tension into a contradiction:
an integer sequence tending to infinity with upward steps of size at most $B$
cannot remain coprime to the pairwise coprime moduli a reduced tail carries.
The
complementary case, in which the error is eventually positive, is the descent
of Theorem~\ref{res:descent}.

\begin{theorem}[bounded negative part]\label{res:bounded}
Let $a,C,D:\N\to\N$ and $E:\N\to\Z$ satisfy
\begin{enumerate}
\item $a_n>1$ and $C_n>0$ for every $n$;
\item the exact dynamics $C_{n+1}+D_n=a_nC_n$ and $D_{n+1}=a_nD_n$;
\item $E_n=\ctr(a_n,D_n,C_n)$ for every $n$;
\item \emph{eventual strict centring}: $|E_n|<C_n$ for all large $n$;
\item \emph{eventually bounded negative part}: $-B\le E_n$ for all large $n$,
for some $B$;
\item \emph{normalised vanishing}: for every $K$ there is an $N$ with
$K\,|E_n|<C_n$ for all $n\ge N$.
\end{enumerate}
Then $E_n=0$ for all sufficiently large $n$.
\end{theorem}

\begin{proof}
Suppose not.  We first remove every late zero: by Theorem~\ref{res:absorb} and
its contrapositive form, $E$ is nowhere zero beyond the centring threshold.
Two cases remain.  If $E$ is negative cofinally,
then~(5) bounds those magnitudes, so Proposition~\ref{res:gcdstab} makes the
tail gcd stabilise and the orbit may be taken reduced past that point.
Hypothesis~(5) with
Proposition~\ref{res:update} gives $C_{n+1}\le C_n+B$, a bounded rise;
hypothesis~(6) with $E$ nowhere zero gives $C_n\to\infty$.  A reduced exact
tail with a divergent numerator and bounded rise is excluded by
Theorem~\ref{res:barrier}.  It remains to treat the case that $E$ is eventually
positive, and there the descent
of Theorem~\ref{res:descent} applies and forces $E$ to vanish, contradicting
nowhere-vanishing.
\end{proof}

\noindent Formalised as the
\lword{Erdos243/ReciprocalTailRigidity.lean}{2265}{boundedNegativePart_eventually_zero}{bounded-negative-part
rigidity}, with the eventual form, in which the centring and the bound need
only hold from some index, as the
\lword{Erdos243/ReciprocalTailRigidity.lean}{2352}{eventuallyBoundedNegativePart_eventually_zero}{eventual
bounded-negative-part rigidity}.  The two auxiliary results it uses are the
\lword{Erdos243/ReciprocalTailRigidity.lean}{1733}{tailState_tendsto_atTop_of_nonzero_normalizedVanishes}{divergence
from normalised vanishing} and the
\lword{Erdos243/ReciprocalTailRigidity.lean}{1748}{no_cofinallyBoundedNegative_of_normalizedVanishes}{exclusion
of cofinally bounded negative magnitudes}, the latter over the
\lword{Erdos243/ReciprocalTailRigidity.lean}{1655}{no_cofinallyBoundedNegative_of_tailTendsto}{tail-divergence
form}.

\begin{corollary}\label{res:cor}
Under the hypotheses of Theorem~\ref{res:bounded}, together with
$C_{n+1}\ne0$ for all large $n$, the multipliers satisfy
$a_{n+1}=a_n^{2}-a_n+1$ for all sufficiently large $n$.
\end{corollary}

\begin{proof}
Theorem~\ref{res:bounded} then Theorem~\ref{res:eventual}.
\end{proof}

Hypotheses~(1)--(3) are the state system, and~(6) is the division-free form of
$|E_n|/C_n\to0$.  Hypothesis~(4) is logically redundant: hypothesis~(6) with
$K=1$ gives $|E_n|<C_n$ eventually.  It is retained because the displayed
statement tracks the linked Lean declaration; shifting to the resulting
threshold gives the same conclusion without an independent centring input.
The formal source does not derive normalised vanishing from
$a_{n+1}/a_n^{2}\to1$, so Theorem~\ref{res:bounded} is, by itself, a
conditional theorem about the state system.  Combined with
Koizumi's bridge below, however, Corollary~\ref{res:cor} settles the
bounded-negative canonical case of Problem~\ref{res:problem}.  What is
unconditional in this note is the all-negative constant and periodic
exclusions, Theorems~\ref{res:constant} and~\ref{res:periodic}, and the barrier
itself, Theorem~\ref{res:barrier}.

The conditionality can nevertheless be located precisely, because the missing
inputs are supplied elsewhere.  With the results of~\cite{koizumi2025},
Theorem~\ref{res:bounded} becomes a conditional theorem about
Problem~\ref{res:problem} itself rather than about the state system alone.  Let
$(a_n)$ satisfy the hypotheses of Problem~\ref{res:problem}.  After deleting a
finite prefix, Corollary~10 of~\cite[p.~8]{koizumi2025} makes the sequence the
pseudo-greedy expansion of its own reciprocal sum, and Lemma~15 there supplies
the integers $c_n,d_n,e_n$ of Section~\ref{sec:problem}.  Hypothesis~(1) then
holds: $c_n$ is a positive integer by Lemma~15(1), and $a_n>1$ after a further
finite shift, because summability of $\sum1/a_n$ forces $a_n\to\infty$.
Hypothesis~(2) is the pair of recurrences $c_{n+1}=a_nc_n-d_n$ and
$d_{n+1}=a_nd_n$ of Lemma~15(2), and hypothesis~(3) is that lemma's
$a_n=(d_n-e_n)/c_n+1$ read as $e_n=d_n-(a_n-1)c_n$, which is the map $\ctr$.
The centring range $-c_n/2\le e_n<c_n/2$ in the same lemma gives
hypothesis~(4) at every index, which is stronger than the eventual form used
here.  Hypothesis~(6) is the vanishing
of the gap sequence in Corollary~10, since $\varepsilon_n=e_n/c_n$.  The sole
hypothesis not supplied is~(5), the eventual bound on the negative part;
nothing in~\cite{koizumi2025} yields it and nothing here does either, so
Problem~\#243 stays open.

\section{Excluding finite normalised negative mass}\label{sec:mass}

Theorem~\ref{res:bounded} uses a uniform bound on the negative part.  A
different finiteness hypothesis, on the total normalised negative mass, removes
a further family of orbits, and the argument is multiplicative rather than
arithmetic: it compares $C$ with a convergent product instead of with a system
of congruences.

\begin{theorem}[finite normalised negative mass]\label{res:mass}
Let \(C_n\in\N_{>0}\) and \(E_n\in\Z\) satisfy
\[
 C_{n+1}=C_n-E_n,\qquad
 \forall K\ \exists N\ \forall n\ge N,\quad K|E_n|<C_n .
\]
If
\[
 \sum_{n=0}^{\infty}\frac{(-E_n)_+}{C_n}<\infty ,
\]
then \(E_n=0\) eventually.  For an exact reciprocal-tail orbit this implies
\(a_{n+1}=a_n^2-a_n+1\) eventually.
\end{theorem}

\begin{proof}
Put \(\delta_n=(-E_n)_+/C_n\).  The update gives
\(C_{n+1}\le C_n(1+\delta_n)\), so
\[
 C_N\le C_0\prod_{n<N}(1+\delta_n).
\]
The product is bounded because \(\sum\delta_n\) converges.  Choose an integer
\(M\) above that bound and apply normalised vanishing with \(K=M\).  Then
\(M|E_n|<C_n<M\) for all sufficiently large \(n\), forcing the integer
\(E_n\) to be zero.  The final recurrence is Theorem~\ref{res:eventual}, whose
hypothesis $C_{n+1}\ne0$ holds because $C_n>0$ throughout.
\end{proof}

\noindent The product argument is the
\lword{Erdos243/SparseResetRecovery.lean}{423}{eventually_zero_of_summable_relativeGrowth}{summable
relative growth}; its specialisation to negative mass is the
\lword{Erdos243/SparseResetRecovery.lean}{488}{eventually_zero_of_summable_negativeRelativeMass}{vanishing
from finite negative mass}, and the statement giving the recurrence directly is
the
\lword{Erdos243/SparseResetRecovery.lean}{510}{sylvesterNext_eventually_of_summable_negativeRelativeMass}{Sylvester
recurrence from summable negative mass}.  This is a checked implication about
the state system, not a derivation of normalised vanishing from
Problem~\#243.  Normalised vanishing is still an input: it is what converts the
bound on $C$ into a bound on $|E_n|$, and nothing here derives it from the
growth condition of Problem~\ref{res:problem}.

The transfer of the previous section applies here as well.  On the canonical
orbit of Corollary~10 of~\cite[p.~8]{koizumi2025} the recurrence
$C_{n+1}=C_n-E_n$ holds with $C_n>0$ by Lemma~15
~\cite[p.~9]{koizumi2025}, and normalised vanishing is the vanishing of
the gap sequence, so Theorem~\ref{res:mass} is a conditional theorem about
Problem~\ref{res:problem} whose sole remaining hypothesis is summability of
$\sum_n(-E_n)_+/C_n$.  That hypothesis is not supplied there, and is the
content of Problem~\ref{res:masshyp} below.

\begin{example}[the product bound at a geometric rate]\label{ex:mass}
Suppose $C_0=100$ and $(-E_n)_+/C_n\le2^{-n-1}$ for every $n$, so that the
normalised negative mass is at most $1$; the constraint at $n=0$ still permits
$E_0$ as negative as $-50$.  Since $1+x\le e^{x}$,
\[
 C_N\le100\prod_{n<N}\bigl(1+2^{-n-1}\bigr)\le100e<272
\]
for every $N$.  Normalised vanishing at $K=272$ then gives
$272|E_n|<C_n<272$ for all large $n$, so $|E_n|<1$ and $E_n=0$ there.  The
constant $272$ comes from the total mass and not from any single magnitude,
which is why the hypothesis of Theorem~\ref{res:mass} is summability of
$(-E_n)_+/C_n$ and not a bound on it.
\end{example}

\section{Complements and further questions}\label{sec:open}

Problem~\#243 is open.  The negative case has narrowed, and it is worth being
precise about what is left.

Each exclusion holds only in its stated regime.  Theorems~\ref{res:constant}
and~\ref{res:periodic} exclude the constant and the stated periodic magnitudes
on an all-negative tail, and do not exclude those patterns merely along the
negative indices of a mixed-sign tail; Theorems~\ref{res:bounded}
and~\ref{res:mass} exclude a bounded negative part and finite normalised
negative mass under normalised vanishing.  Together with absorption and descent
they fix the profile of any counterexample.

\begin{proposition}[frontier profile for a counterexample]\label{res:frontier}
For the canonical integer state attached to any counterexample to
Problem~\ref{res:problem},
\[
 E_n\ne0\quad\hbox{eventually},\qquad
 \frac{|E_n|}{C_n}\longrightarrow0,
\]
and
\[
 \limsup_{\substack{n\to\infty\\E_n<0}}(-E_n)=\infty,
 \qquad
 \sum_{n=0}^{\infty}\frac{(-E_n)_+}{C_n}=\infty .
\]
In particular, negative indices occur infinitely often.  Thus the remaining
regime consists of globally coherent, unbounded negative excursions with
divergent normalised negative mass.
\end{proposition}

\begin{proof}
Koizumi's canonical-tail transfer gives strict centring and normalised
vanishing after a finite shift; the strict-centring hypothesis of
Theorem~\ref{res:bounded} also follows from normalised vanishing with $K=1$.
Absorption then makes $E_n\ne0$ eventually,
since eventual zero would give the Sylvester recurrence.  Descent excludes an
eventually nonnegative state.  Theorem~\ref{res:bounded} excludes a bounded
negative part, and Theorem~\ref{res:mass} excludes finite normalised negative
mass.  These statements are unchanged by deleting a finite prefix.
\end{proof}

Put
\[
 \delta_n:=\frac{(-E_n)_+}{C_n}.
\]
Every argument in Sections~\ref{sec:constant}--\ref{sec:bounded} uses a
finiteness hypothesis: a fixed set of prime divisors, a fixed period, or a
fixed bound on the negative part.  Theorem~\ref{res:mass} uses
$\sum_n\delta_n<\infty$.  Proposition~\ref{res:frontier} records that a
counterexample has none of them; it is a frontier statement, not an additional
problem equivalent to the original one.  Stated at the level of state orbits,
the surviving obstruction is the following.

\begin{problem}[unbounded divergent-mass negative excursions]\label{res:excursions}
Exclude, or construct, an exact natural orbit $(a,D,C)$ with $a_n>1$ and
$C_n>0$, whose multipliers satisfy $\lim_n a_{n+1}/a_n^{2}=1$ and whose centred
state satisfies normalised vanishing, and which is
negative infinitely often with magnitudes unbounded along that cofinal set and
\[
 \sum_n\frac{(-E_n)_+}{C_n}=\infty .
\]
An excursion of this kind has none of the finiteness hypotheses just listed.
By Proposition~\ref{res:frontier} the canonical state of a counterexample is
such an orbit after deletion of a finite prefix, so an exclusion would settle
Problem~\#243.  A construction would not settle it in the other direction: an
exact state orbit need not arise from a sequence satisfying
Problem~\ref{res:problem}.
\end{problem}

\begin{remark}
The hypotheses in Problem~\ref{res:excursions} beyond the state recurrence are
not decorative: the recurrence alone admits unbounded divergent-mass
excursions.  Put
\[
 C_n=2^{n},\qquad b_0=2,\qquad b_{n+1}=\tfrac12 b_n(b_n+2),\qquad
 a_n=b_n+2,\qquad D_n=b_nC_n ,
\]
so that $(b_n)=2,4,12,84,3612,\ldots$ and $(a_n)=4,6,14,86,3614,\ldots$.  Every
$b_n$ is a positive even integer, since $b_n=2k$ gives $b_{n+1}=2k(k+1)$, so all
four sequences are $\N$-valued.  The orbit is exact:
$C_{n+1}+D_n=2^{n}(2+b_n)=a_nC_n$ and
$D_{n+1}=b_{n+1}2^{n+1}=b_n(b_n+2)2^{n}=a_nD_n$.  Its centred state is
\[
 E_n=D_n-(a_n-1)C_n=b_n2^{n}-(b_n+1)2^{n}=-C_n ,
\]
so the negative magnitudes $2^{n}$ are unbounded and $\sum_n(-E_n)_+/C_n$
diverges.  What fails is the centring: $|E_n|=C_n$ at every index, so strict
centring fails at the boundary and normalised vanishing fails outright.  The
multipliers satisfy $a_{n+1}=\tfrac12(a_n^{2}-2a_n+4)$, so
$a_{n+1}/a_n^{2}\to\tfrac12$ and the growth hypothesis fails as well.  The
example therefore says nothing about Problem~\ref{res:problem}, and it does not
isolate the centring hypotheses: $E_n=-C_n$ at every index forces that
multiplier recurrence, so growth fails with them.  It is recorded only to show
that the state recurrence alone is not enough.
\end{remark}

\subsection{Global height and old-factor overlap}

The main arithmetic gap is global.  Define the cumulative digit LCM and its
overlap quotient by
\[
 L_0=D_0,\qquad L_{n+1}=\lcm(L_n,a_n),\qquad M_n=\frac{D_n}{L_n}.
\]
Since $D_n=D_0\prod_{j<n}a_j$, the quotient is integral and
$M_nL_n=D_n$.  It records the multiplicity lost when the full product scale
is compressed to an LCM.

\subsection{Repair entropy across a recovery}

The cumulative quotient does not by itself distinguish one large deletion
from many small ones.  The reduced recurrence does.  Let $u,h,c:\N\to\N$ and
$e:\N\to\Z$ satisfy
\[
 h_nu_{n+1}=u_n-e_n,
 \qquad u_n>0,
 \qquad h_n>0.
\]
Fix $K,r,L\in\N$ with $K,L>0$, assume
\[
 K|e_{r+i}|<u_{r+i}\quad(0\le i<L),
 \qquad u_r\le u_{r+L},
\]
and let $R$ be a finite family of moduli.  Suppose every $m_q$, $q\in R$,
divides the deletion product $\prod_{r\le n<r+L}c_n$, and that
$c_n^2\mid h_n$ throughout the interval.  Then
\begin{equation}\label{eq:repair-entropy}
 K^L\lcm(m_q:q\in R)^2<(K+1)^L.
\end{equation}
This is the
\lword{Erdos243/RepairEntropy.lean}{48}
      {repairedFamily_recovery_energy_divisionFree}
      {repair-entropy conservation law}.

The square in~\eqref{eq:repair-entropy} is the key point.  Exact denominator
reduction charges $c_n^2$ to the scale growth, so a repaired family whose LCM
is at least $2^{|R|}$ satisfies
\[
 K^L4^{|R|}<(K+1)^L.
\]
The formal consequence is
\lword{Erdos243/RepairEntropy.lean}{78}
      {repaired_card_bound_of_independent}
      {the independent-repair cardinality bound}.
It rules out positive-density deletion of independent old moduli inside
recoveries with a fixed relative-error budget.

There is also an endpoint that does not mention a chosen modulus family.  If
$K|e_n|<u_n$ eventually for every $K$, then for each fixed $L>0$ there is an
$N$ such that every recovery $u_r\le u_{r+L}$ with $r\ge N$ satisfies
\[
 \prod_{0\le i<L}h_{r+i}=1.
\]
Thus a surviving nontrivial reset must have recovery length tending to
infinity.  Lean checks this as
\lword{Erdos243/RepairEntropy.lean}{107}
      {eventually_recoveryPayment_eq_one_of_fixedLength}
      {eventual disappearance of fixed-length reset payments}.
The theorem does not bound a recovery length that varies with $r$, nor does it
supply a lower bound for the LCM of the moduli repaired in such an interval.
Those are the missing inputs needed to turn~\eqref{eq:repair-entropy} into a
global contradiction.

\begin{problem}[global overlap-height growth]\label{res:lcmheight}
For every nonterminal canonical orbit satisfying Problem~\ref{res:problem},
must
\[
 \limsup_{n\to\infty}\frac{\log M_n}{n}>0?
\]
Equivalently, must there be a $K\ge1$ for which
\[
 2^n\le M_n^K
\]
at infinitely many indices?
\end{problem}

\noindent The two displayed formulations are equivalent up to changing the
positive constant; the second is not a weaker target. On the canonical orbit,
Section~\ref{sec:lcmrecords} supplies $M_n\mid C_n$ and $M_n\le C_n$.
Therefore a positive answer would contradict the subexponential tail-height
budget and prove Problem~\#243. The missing input is the proposed lower
bound on overlap growth.

There is also a more local-looking question whose content is nevertheless the
entire prefix.  Put $A_n=\prod_{j<n}a_j$, so $D_n=D_0A_n$.  From
\[
 D_0A_n=(a_n-1)C_n+E_n
\]
one immediately obtains
\[
 \gcd(A_n,a_n-1)\mid E_n.
\]
The checked
\lword{Erdos243/ReciprocalTailRigidity.lean}{2164}{tailState_subexponential_of_normalizedVanishes}{tail-height estimate}
and normalised vanishing make $|E_n|$ subexponential in $n$.

\begin{problem}[old-factor overlap]\label{res:prefixgcd}
In every nonterminal canonical orbit satisfying Problem~\ref{res:problem}, is
\[
 \limsup_{n\to\infty}
 \frac{\log\gcd(A_n,a_n-1)}{n}>0?
\]
Equivalently, do there exist $\eta>0$ and infinitely many $n$ such that
$\gcd(A_n,a_n-1)\ge e^{\eta n}$?
\end{problem}

\noindent A positive answer contradicts the displayed divisibility and the
subexponential error budget.  Arbitrarily long locally admissible blocks do
not answer this question: the gcd uses the complete prefix.

\subsection{The direct analytic question}

\begin{problem}[summability from rationality]\label{res:masshyp}
Let $(a_n)$ satisfy Problem~\ref{res:problem}, and let $(D_n,C_n,E_n)$ be its
canonical integer state.  Must
\[
 \sum_{n=0}^{\infty}\delta_n
 =\sum_{n=0}^{\infty}\frac{(-E_n)_+}{C_n}<\infty ?
\]
\end{problem}

\noindent An affirmative answer closes Problem~\#243 by
Theorem~\ref{res:mass}.  Since $\delta_n\to0$, convergence is equivalently
expressible as boundedness of the partial products
$\prod_{n<N}(1+\delta_n)$; this is the same criterion in multiplicative form.
A negative answer requires a sequence satisfying the full growth and
rationality hypotheses. A locally admissible state orbit is insufficient.
By Theorem~\ref{res:weightedrecord}, it is enough to establish the finiteness
of~\eqref{eq:weightedgrowth} for one decreasing weight with divergent
integral and one fixed baseline. This replaces the all-index mass by a
discounted series supported only at LCM records.

\begin{remark}
One further question is not left open here, since it is answered
in~\cite{koizumi2025}: whether normalised vanishing follows from
$a_{n+1}/a_n^{2}\to1$ and rationality.  It does.  After
deleting a finite prefix, Koizumi's Corollary~10
~\cite[p.~8]{koizumi2025} makes the sequence the pseudo-greedy expansion of its
own reciprocal sum and gives $\varepsilon_n\to0$ under exactly the hypotheses of
Problem~\ref{res:problem}; rationality is not needed for that step, only
summability of $\sum1/a_n$.  Since $\varepsilon_n=E_n/C_n$ under the intended
reading of Section~\ref{sec:state}, this is normalised vanishing.  Strict
centring comes from the same source and is stronger than what
Theorem~\ref{res:bounded} assumes: for a rational reciprocal sum, the range
$-c_n/2\le e_n<c_n/2$ of his Lemma~4(2)
~\cite[p.~11]{koizumi2025} gives $|E_n|<C_n$ at every index, which is stronger
than eventual strict centring.  The eventual form also follows directly from
normalised vanishing with $K=1$.  What remains unsupplied is hypothesis~(5), the bound
on the negative part, which is why the theorem is still conditional.
\end{remark}

\subsection{A barrier extension}

\begin{problem}[variable-rise CRT barrier]\label{res:variablerise}
Determine the weakest growth condition on the positive increments of $u_n$
under which an unbounded integer sequence cannot remain coprime to infinitely
many pairwise-coprime old moduli as in Theorem~\ref{res:barrier}.  In
particular, does the conclusion remain valid under the candidate condition
\[
 (u_{n+1}-u_n)_+=o\bigl(\log\log(u_n+3)\bigr)?
\]
\end{problem}

\noindent The displayed rate is a proposed threshold, not a proved sharp
boundary.  Either an extension of the CRT barrier or a counterexample with
this rise rate would be informative independently of Problem~\#243.

\subsection*{Formalisation targets}

\paragraph{Erd\H{o}s--Straus.}
Formalise Theorem~3 of~\cite{erdosstraus1964}, with its prefix-LCM quotient
and correctly indexed growth factor, under the published analytic hypotheses.

\paragraph{Duverney.}
Formalise Corollary~3.2 of~\cite{duverney2001}, including its signed
numerators and its one-sided summability hypothesis, and then recover the
all-positive specialisation used here.  These are useful verification
projects, but neither is an additional open mathematical problem.

\pdfbookmark[1]{Statements and declarations}{statements-declarations}
\subsection*{Statements and declarations}

\paragraph{Artefact and data availability.} The
\href{\sourceurl}{pinned formal-source revision} contains the Lean sources, the
fixed toolchain, and the library manifest used in the verification.  Proof
authority remains with those sources; this manuscript provides navigation and
exposition.

Lean checks each proof term against the fixed library
version, and the sources linked here contain no proof placeholders and no
project-defined axioms; Lean does not authorise the exposition, the citation
choices, or the interpretation, for which the author remains responsible.

\paragraph{Funding and competing interests.} This work received no external
funding.  The author declares no competing interests.

\paragraph{Acknowledgements.} The problem numbering and status follow the
Erd\H{o}s Problems catalogue maintained by Thomas Bloom~\cite{erdosproblems}.

\appendix
\section{Guide to the formal sources}\label{app:index}

Each linked phrase opens its Lean declaration at the pinned source
revision~\commitshort.  The state system, the exclusions, the barrier, and the
bounded-negative-part theorem are in \texttt{ReciprocalTailRigidity.lean}; the
finite-negative-mass theorem is in \texttt{SparseResetRecovery.lean}; and the
repair-entropy inequalities are in \texttt{RepairEntropy.lean}.  The residue
reduction of Appendix~\ref{app:residue} is in
\texttt{FiniteHorizonResidue.lean}.  Three distinctions are worth carrying into
the source.  The
identification of any of these modules with reciprocal tails is the exposition
of Section~\ref{sec:state} and is not a checked statement.  The periodic exclusion
assumes the regime $e_n<a_n$.  And the final Lean declaration lists strict
centring and normalised vanishing as hypotheses; the former follows eventually
from the latter with $K=1$, while the formal source derives neither from the
original analytic problem.  For the external bridge supplying them,
see Sections~\ref{sec:problem} and~\ref{sec:bounded}.

% The exhaustive source manifest is retained in the authored TeX for automated
% validation, and kept outside the printed note.
\iffalse

\begin{tabular}{@{}ll@{}}
informal statement & linked Lean source\\
Sylvester successor & \lrefx{Erdos243/ReciprocalTailRigidity.lean}{37}{sylvesterNext}\\
denominator update & \lrefx{Erdos243/ReciprocalTailRigidity.lean}{41}{nextDenState}\\
tail update & \lrefx{Erdos243/ReciprocalTailRigidity.lean}{45}{nextTailState}\\
centred state & \lrefx{Erdos243/ReciprocalTailRigidity.lean}{49}{centeredState}\\
Sylvester defect & \lrefx{Erdos243/ReciprocalTailRigidity.lean}{53}{sylvesterDefect}\\
update law & \lrefx{Erdos243/ReciprocalTailRigidity.lean}{57}{nextTailState_eq_sub_centered}\\
denominator equivariance & \lrefx{Erdos243/ReciprocalTailRigidity.lean}{65}{nextDenState_scale}\\
tail equivariance & \lrefx{Erdos243/ReciprocalTailRigidity.lean}{70}{nextTailState_scale}\\
centred equivariance & \lrefx{Erdos243/ReciprocalTailRigidity.lean}{75}{centeredState_scale}\\
normalised exclusion & \lrefx{Erdos243/ReciprocalTailRigidity.lean}{100}{no_normalizedConstantNegative_orbit}\\
scale-primitive exclusion & \lrefx{Erdos243/ReciprocalTailRigidity.lean}{147}{no_scalePrimitiveConstantNegative_orbit}\\
constant exclusion at every scale & \lrefx{Erdos243/ReciprocalTailRigidity.lean}{286}{no_constantNegative_orbit}\\
eventual constant exclusion & \lrefx{Erdos243/ReciprocalTailRigidity.lean}{350}{no_eventuallyConstantNegative_orbit}\\
persistence of a multiplier & \lrefx{Erdos243/ReciprocalTailRigidity.lean}{382}{base_dvd_denState_later}\\
one-step lock & \lrefx{Erdos243/ReciprocalTailRigidity.lean}{402}{divisorLocksTailSucc}\\
persistence of a lock & \lrefx{Erdos243/ReciprocalTailRigidity.lean}{429}{divisorLock_persists}\\
iterated period & \lrefx{Erdos243/ReciprocalTailRigidity.lean}{461}{periodic_add_mul}\\
iterated drift & \lrefx{Erdos243/ReciprocalTailRigidity.lean}{476}{phaseDrift_add_mul}\\
repeated-divisor transport & \lrefx{Erdos243/ReciprocalTailRigidity.lean}{494}{repeatedDivisor_forces_periodicCommonScale}\\
phase-primitive exclusion & \lrefx{Erdos243/ReciprocalTailRigidity.lean}{548}{no_phasePrimitivePeriodicNegative_orbit}\\
periodic exclusion at every scale & \lrefx{Erdos243/ReciprocalTailRigidity.lean}{690}{no_periodicNegative_orbit}\\
eventual periodic exclusion & \lrefx{Erdos243/ReciprocalTailRigidity.lean}{796}{no_eventuallyPeriodicNegative_orbit}\\
shifted consecutive multiples & \lrefx{Erdos243/ReciprocalTailRigidity.lean}{833}{exists_shifted_consecutiveMultiples}\\
bounded-rise barrier & \lrefx{Erdos243/ReciprocalTailRigidity.lean}{897}{no_boundedRise_of_tailAvoidance}\\
step coprimality & \lrefx{Erdos243/ReciprocalTailRigidity.lean}{983}{reducedStep_coprime_currentFactor}\\
pairwise coprimality & \lrefx{Erdos243/ReciprocalTailRigidity.lean}{1010}{reducedTail_pairwiseCoprime}\\
whole-modulus avoidance & \lrefx{Erdos243/ReciprocalTailRigidity.lean}{1024}{reducedTail_wholeModulusAvoidance}\\
reduced bounded-rise exclusion & \lrefx{Erdos243/ReciprocalTailRigidity.lean}{1035}{no_boundedRise_reducedTail}\\
eventual reduced exclusion & \lrefx{Erdos243/ReciprocalTailRigidity.lean}{1060}{no_eventuallyBoundedRise_reducedTail}\\
gcd monotonicity & \lrefx{Erdos243/ReciprocalTailRigidity.lean}{1424}{tailGcd_dvd_succ}\\
gcd at a negative index & \lrefx{Erdos243/ReciprocalTailRigidity.lean}{1581}{tailGcd_eq_gcd_negativeMagnitude}\\
chain stabilisation & \lrefx{Erdos243/ReciprocalTailRigidity.lean}{1599}{dvdChain_eventuallyConstant_of_cofinally_bounded}\\
gcd stabilisation & \lrefx{Erdos243/ReciprocalTailRigidity.lean}{1630}{tailGcd_eventuallyConstant_of_cofinally_boundedNegative}\\
tail-divergence form & \lrefx{Erdos243/ReciprocalTailRigidity.lean}{1655}{no_cofinallyBoundedNegative_of_tailTendsto}\\
divergence from normalised vanishing & \lrefx{Erdos243/ReciprocalTailRigidity.lean}{1733}{tailState_tendsto_atTop_of_nonzero_normalizedVanishes}\\
exclusion from normalised vanishing & \lrefx{Erdos243/ReciprocalTailRigidity.lean}{1748}{no_cofinallyBoundedNegative_of_normalizedVanishes}\\
defect identity & \lrefx{Erdos243/ReciprocalTailRigidity.lean}{1769}{sylvesterDefect_mul_nextTailState}\\
local rigidity step & \lrefx{Erdos243/ReciprocalTailRigidity.lean}{1781}{sylvesterNext_eq_of_centered_zero}\\
eventual Sylvester recurrence & \lrefx{Erdos243/ReciprocalTailRigidity.lean}{1799}{sylvesterNext_eventually_of_centered_zero}\\
eventual constancy & \lrefx{Erdos243/ReciprocalTailRigidity.lean}{1819}{antitone_nat_eventually_constant}\\
stabilisation of the nonnegative error & \lrefx{Erdos243/ReciprocalTailRigidity.lean}{1837}{centeredState_eventually_zero}\\
tail realisation & \lrefx{Erdos243/ReciprocalTailRigidity.lean}{1863}{natTail_eq_nextTailState}\\
denominator realisation & \lrefx{Erdos243/ReciprocalTailRigidity.lean}{1875}{natDen_eq_nextDenState}\\
signed update law & \lrefx{Erdos243/ReciprocalTailRigidity.lean}{1972}{natTail_eq_sub_centeredState}\\
absorption of a vanishing centred state & \lrefx{Erdos243/ReciprocalTailRigidity.lean}{2221}{centeredState_zero_absorbing}\\
contrapositive along a tail & \lrefx{Erdos243/ReciprocalTailRigidity.lean}{2248}{eventually_nonzero_of_zero_absorbing}\\
bounded-negative-part rigidity & \lrefx{Erdos243/ReciprocalTailRigidity.lean}{2265}{boundedNegativePart_eventually_zero}\\
eventual bounded-negative-part rigidity & \lrefx{Erdos243/ReciprocalTailRigidity.lean}{2352}{eventuallyBoundedNegativePart_eventually_zero}\\
repair-entropy conservation & \lrefx{Erdos243/RepairEntropy.lean}{48}{repairedFamily_recovery_energy_divisionFree}\\
independent-repair bound & \lrefx{Erdos243/RepairEntropy.lean}{78}{repaired_card_bound_of_independent}\\
fixed-length reset disappearance & \lrefx{Erdos243/RepairEntropy.lean}{107}{eventually_recoveryPayment_eq_one_of_fixedLength}\\
forced numerator & \lrefx{Erdos243/FiniteHorizonResidue.lean}{22}{forcedNumerator}\\
polynomial congruence & \lrefx{Erdos243/FiniteHorizonResidue.lean}{26}{forcedNumerator_modEq}\\
exact-division cancellation & \lrefx{Erdos243/FiniteHorizonResidue.lean}{41}{quotient_modEq_of_modEq_mul}\\
horizon modulus & \lrefx{Erdos243/FiniteHorizonResidue.lean}{53}{horizonModulus}\\
ascending-factorial form & \lrefx{Erdos243/FiniteHorizonResidue.lean}{59}{horizonModulus_eq_ascFactorial}\\
factorial value at the initial index & \lrefx{Erdos243/FiniteHorizonResidue.lean}{69}{horizonModulus_zero_eq_factorial}\\
survival predicate & \lrefx{Erdos243/FiniteHorizonResidue.lean}{75}{ForcedSurvives}\\
quotient transport & \lrefx{Erdos243/FiniteHorizonResidue.lean}{85}{forcedQuotient_modEq}\\
shrinking-modulus transport & \lrefx{Erdos243/FiniteHorizonResidue.lean}{97}{forcedSurvives_iff_of_modEq}\\
factorial residue reduction & \lrefx{Erdos243/FiniteHorizonResidue.lean}{134}{forcedSurvives_iff_of_modEq_factorial}\\
\end{tabular}
\fi

\section{A factorial residue reduction for forced orbits}\label{app:residue}

In the case $m=c=1$ of Section~\ref{sec:constant}, where the shape
equation~\eqref{eq:shape} reads $D_n+1=(a_n-1)(n+1)$, each multiplier is
determined by its predecessor; we call such an orbit \emph{forced}.  At index
$n$ the numerator of the next multiplier is
\[
 \fnum(n,a)=(n+1)a^{2}-(n+2)a+(n+3),
\]
the
\lword{Erdos243/FiniteHorizonResidue.lean}{22}{forcedNumerator}{forced
numerator}, and the divisor is $n+2$.  The orbit survives a step when that
division is exact, giving a survival predicate, the
\lword{Erdos243/FiniteHorizonResidue.lean}{75}{ForcedSurvives}{survival
predicate}.  Example~\ref{ex:shape} is the forced orbit from $a=3$ read this
way: $\fnum(0,3)=6$ is divisible by $2$ and gives $a_1=3$, while
$\fnum(1,3)=13$ is not divisible by $3$, so the orbit stops there.  Deciding
survival by iteration is expensive because the orbit
grows doubly exponentially, and it is unnecessary: survival over a finite
horizon depends on the initial value only through a factorial residue.

\begin{theorem}[factorial residue reduction]\label{res:residue}
For all $h$ and all integers $a\equiv b \pmod{(h+1)!}$, the orbit from $a$
survives $h$ forced updates exactly when the orbit from $b$ does.
\end{theorem}

\begin{proof}
Let $M(0,i)=1$ and $M(h+1,i)=(i+2)M(h,i+1)$, an ascending factorial with
$M(h,0)=(h+1)!$.  The numerator is a polynomial with integer coefficients, so
congruences transfer; reducing the modulus gives divisibility by $i+2$ for one
exactly when for the other, and cancelling that common factor from both values
and modulus leaves the inductive hypothesis at $M(h,i+1)$.
\end{proof}

\noindent Formalised as the
\lword{Erdos243/FiniteHorizonResidue.lean}{134}{forcedSurvives_iff_of_modEq_factorial}{factorial
residue reduction}, over the
\lword{Erdos243/FiniteHorizonResidue.lean}{97}{forcedSurvives_iff_of_modEq}{shrinking-modulus
transport}, the
\lword{Erdos243/FiniteHorizonResidue.lean}{26}{forcedNumerator_modEq}{polynomial
congruence}, and the
\lword{Erdos243/FiniteHorizonResidue.lean}{41}{quotient_modEq_of_modEq_mul}{exact-division
cancellation}; the modulus identifications are the
\lword{Erdos243/FiniteHorizonResidue.lean}{59}{horizonModulus_eq_ascFactorial}{ascending-factorial
form} and the
\lword{Erdos243/FiniteHorizonResidue.lean}{69}{horizonModulus_zero_eq_factorial}{factorial
value at the initial index}.

At $h=1$ the modulus is $2!=2$, and surviving one update means
$2\mid a^{2}-2a+3$, which holds exactly for odd $a$: for instance
$\fnum(0,3)=6$ but $\fnum(0,4)=11$.  So one step of survival is decided by the
parity of $a$ alone, which is Theorem~\ref{res:residue} at its smallest
nontrivial horizon.

The search this supported is superseded.  Running it over initial states below
$5000$ produced forced prefixes of length $17$ and no longer;
Theorem~\ref{res:constant} now excludes the
constant-negative case outright, for every seed and at every scale.  The
reduction is retained because it is exact, and because the shrinking-modulus
technique transfers to any forced orbit whose step is a polynomial division.

% ---- part family_catalogue ----
\clearpage
\section{Complete result-family map}\label{sec:erdos-243-complete-family-map}

This section places every registered family for this problem in the shared 93-family reader order.  The five display bands control exposition only; the separate promotion state currently covers 23 families and is reported but does not hide or strengthen any family.  Mathematical statements and evidence modes come from the public claim registry.  Across all eight problems the public result-atom catalog contains 704 exact packet coordinates; this problem contributes 84.  Catalog rows expose bounded statement excerpts plus full-source digests, not a claim that every complete packet statement is reproduced here.

\subsection{Centered state dynamics}
\textbf{Reader position.} 2 of 93; display band: front door.  Formal editorial disposition: split.  These are separate classifications.

\textbf{Reader entry.} Track one signed centered error: forcing it eventually to zero recovers the Sylvester recurrence.

Under exact natural C/D dynamics with a > 1, C > 0, E = D - (a - 1)C, strict |E| < C, a uniform lower bound on E, and normalized vanishing, boundedNegativePart\_eventually\_zero forces E = 0 eventually. centeredState\_eventually\_zero supplies the nonnegative stabilization mechanism, and sylvesterNext\_eventually\_of\_centered\_zero recovers eventual Sylvester recurrence when the next tail is nonzero.

\textbf{Authority and reach.} Lean kernel; Comparator-selected; locally proved signed recovery/rigidity theorem; novelty unassessed.

\textbf{Exact boundary.} The signed theorem has a C D : \ensuremath{\mathbb{N}} \ensuremath{\to} \ensuremath{\mathbb{N}}, E : \ensuremath{\mathbb{N}} \ensuremath{\to} \ensuremath{\mathbb{Z}}, and B : \ensuremath{\mathbb{N}} with all exact hypotheses explicit: a > 1, C > 0, C(n + 1) + D(n) = a(n) \ensuremath{\mathbin{\cdot}} C(n), D(n + 1) = a(n) \ensuremath{\mathbin{\cdot}} D(n), E is the centered state D \ensuremath{-} (a \ensuremath{-} 1)C, |E| < C, a uniform lower bound on E, and division-free normalized vanishing. It excludes persistent bounded negative behaviour and forces eventual centered defect zero; centered-zero plus an eventually nonzero next tail recovers Sylvester, but no reciprocal-tail irrationality follows. The unbounded negative mixed-sign branch and the prime-specific producer remain open. This is distinct from the periodic negative-orbit and cofinally bounded-negative no-go families; no novelty, priority, significance, or external-review claim is made, and no unrestricted \#243 solution follows.

\textbf{Result-atom population.} 32 of 704 public coordinates. The public atom catalog groups them under this family.

\textbf{Comparator assurance.} exact selected interface; 1 executable interface(s), 1 repository-registered selected result interface(s).

\textbf{Publication placement.} This family is also admitted to the short note.

\begin{itemize}
  \item \nolinkurl{ErdosProblems.Erdos243.boundedNegativePart_eventually_zero}
  \item \nolinkurl{ErdosProblems.Erdos243.centeredState_eventually_zero}
  \item \nolinkurl{ErdosProblems.Erdos243.sylvesterNext_eventually_of_centered_zero}
\end{itemize}

\subsection{Negative mass recovery}
\textbf{Reader position.} 11 of 93; display band: major result.  Formal editorial disposition: promote.  These are separate classifications.

\textbf{Reader entry.} If the normalized negative centered mass is summable, the negative excursions must end and the Sylvester recurrence follows.

Under exact reciprocal-tail C/D dynamics, positive tail, the strict centered-step equation, and division-free normalized vanishing, summability of the normalized negative centered mass forces centeredState to vanish eventually and recovers the Sylvester recurrence. The load-bearing theorem is supported by the exact tail-growth bound and the finite-negative-mass stabilization theorem.

\textbf{Authority and reach.} Lean kernel; conditional analytic recovery criterion; novelty unassessed.

\textbf{Exact boundary.} The exact criterion assumes a, D : \ensuremath{\mathbb{N}} \ensuremath{\to} \ensuremath{\mathbb{Z}} and C : \ensuremath{\mathbb{N}} \ensuremath{\to} \ensuremath{\mathbb{N}} with the declared next-denominator and next-tail dynamics, C > 0, the strict centered step equation, division-free normalized vanishing, and Summable (negativeRelativeMass C (fun n \ensuremath{\mapsto} centeredState (a n) (D n) (C n))). It does not prove these hypotheses for the original Erdős \#243 orbit, supply a prime-specific producer, or establish an unconditional \#243 endpoint or reciprocal-tail irrationality. Any surviving canonical orbit must have divergent normalized negative mass. This is distinct from centered\_state\_dynamics, whose recovery criterion uses a uniform lower-bound hypothesis; no novelty, priority, significance, or external-review claim is made.

\textbf{Result-atom population.} 36 of 704 public coordinates. The public atom catalog groups them under this family.

\textbf{Comparator assurance.} exact selected interface; 1 executable interface(s), 1 repository-registered selected result interface(s).

\textbf{Publication placement.} This family is also admitted to the short note.

\begin{itemize}
  \item \nolinkurl{ErdosProblems.Erdos243.tail_growth_le_one_add_negativeRelativeMass}
  \item \nolinkurl{ErdosProblems.Erdos243.eventually_zero_of_summable_negativeRelativeMass}
  \item \nolinkurl{ErdosProblems.Erdos243.sylvesterNext_eventually_of_summable_negativeRelativeMass}
\end{itemize}

\subsection{Bounded negative exclusion}
\textbf{Reader position.} 15 of 93; display band: major result.  Formal editorial disposition: promote.  These are separate classifications.

\textbf{Reader entry.} A negative centered error cannot stay bounded forever under the stated normalized dynamics.

Normalised vanishing excludes a cofinally bounded negative part and forces eventual Sylvester behaviour.

\textbf{Authority and reach.} Lean kernel; Comparator-selected; conditional reduction and no-go result.

\textbf{Exact boundary.} Every positivity, dynamics, bounded-rise, and vanishing hypothesis remains explicit.

\textbf{Result-atom population.} 2 of 704 public coordinates. The public atom catalog groups them under this family.

\textbf{Comparator assurance.} exact selected interface; 1 executable interface(s), 1 repository-registered selected result interface(s).

\textbf{Publication placement.} This family is also admitted to the short note.

\subsection{Arithmetic weighted record dichotomy}
\textbf{Reader position.} 21 of 93; display band: major result.  Formal editorial disposition: promote.  These are separate classifications.

\textbf{Reader entry.} For positive integer LCM states with exact lcm/gcd update and feedback rho U\_next=U-V, V=L-(a-1)U, and only the lower centering bound -U<=2V, bounded U is equivalent to finiteness of the record-step excess sum sum(-V-B)\_+ f(U), for each fixed integer B>=0 and finite nonnegative nonincreasing weight with divergent integral. Unbounded height itself supplies arbitrarily many pairwise-coprime record digits; no fresh-index density or subexponential growth assumption is needed.

For positive integer LCM states with exact lcm/gcd update and feedback rho U\_next=U-V, V=L-(a-1)U, and only the lower centering bound -U<=2V, bounded U is equivalent to finiteness of the record-step excess sum sum(-V-B)\_+ f(U), for each fixed integer B>=0 and finite nonnegative nonincreasing weight with divergent integral. Unbounded height itself supplies arbitrarily many pairwise-coprime record digits; no fresh-index density or subexponential growth assumption is needed.

\textbf{Authority and reach.} Ordinary proof with independent exact finite checks; ordinary arithmetic theorem and canonical-orbit criterion.

\textbf{Exact boundary.} The finite CRT crossing and global divergent-integral argument are stated in the public paper. The full dichotomy has no Lean or external challenge certificate. Turning the canonical weighted criterion into universal Sylvester termination still requires a genuine upper summability bound.

\textbf{Result-atom population.} 1 of 704 public coordinates. The public atom catalog groups them under this family.

\textbf{Comparator assurance.} not applicable to comparator; 0 executable interface(s), 0 repository-registered selected result interface(s).

\textbf{Publication placement.} This family remains in the complete long record and is not a short-note headline.

\subsection{Bounded rise coprimality}
\textbf{Reader position.} 43 of 93; display band: mechanism.  Formal editorial disposition: merge.  These are separate classifications.

\textbf{Reader entry.} Fresh coprime moduli eventually defeat a reduced tail that rises too slowly.

A bounded-rise sequence cannot remain coprime to fresh pairwise-coprime moduli; reduced tails inherit this obstruction.

\textbf{Authority and reach.} Lean kernel; locally proved result; novelty unassessed.

\textbf{Exact boundary.} The required bounded-rise input is not automatic for the original sequence.

\textbf{Result-atom population.} 11 of 704 public coordinates. The public atom catalog groups them under this family.

\textbf{Comparator assurance.} represented by selected interface; 0 executable interface(s), 0 repository-registered selected result interface(s).

\textbf{Publication placement.} This family is also admitted to the short note.

\begin{itemize}
  \item \nolinkurl{ErdosProblems.Erdos243.no_eventuallyBoundedRise_reducedTail}
  \item \nolinkurl{ErdosProblems.Erdos243.reducedTail_pairwiseCoprime}
\end{itemize}

\subsection{Negative orbit no go}
\textbf{Reader position.} 70 of 93; display band: frontier.  Formal editorial disposition: promote.  These are separate classifications.

\textbf{Reader entry.} Periodic negative-error patterns cannot be the obstruction sought in Erdős 243.

Constant, eventually constant, periodic, and eventually periodic negative-error orbits are excluded in their stated regimes.

\textbf{Authority and reach.} Lean kernel; Comparator-selected; no-go result.

\textbf{Exact boundary.} The exclusions do not cover arbitrary unbounded negative behaviour.

\textbf{Result-atom population.} 2 of 704 public coordinates. The public atom catalog groups them under this family.

\textbf{Comparator assurance.} exact selected interface; 1 executable interface(s), 1 repository-registered selected result interface(s).

\textbf{Publication placement.} This family is also admitted to the short note.

% ---- part back ----
\begin{thebibliography}{9}
\small
\raggedright
\setlength{\itemsep}{0pt}
\bibitem{erdosstraus1964} P.~Erd\H{o}s and E.~G. Straus,
  \href{https://users.renyi.hu/~p_erdos/1964-19.pdf}{\emph{On the irrationality
  of certain Ahmes series}},
  J. Indian Math. Soc. (N.S.) \textbf{27} (1964), 129--133. MR~175848.
\bibitem{duverney2001} D.~Duverney,
  \href{https://www.ms.u-tokyo.ac.jp/journal/pdf/jms080206.pdf}
  {\emph{Irrationality of fast converging series of rational numbers}},
  J. Math. Sci. Univ. Tokyo \textbf{8} (2001), 275--316. MR~1837165.
\bibitem{badea1993} C.~Badea,
  \href{https://matwbn.icm.edu.pl/ksiazki/aa/aa63/aa6342.pdf}{\emph{A theorem
  on irrationality of infinite series and applications}}, Acta Arith.
  \textbf{63} (1993), no.~4, 313--323.
\bibitem{erdosgraham1980} P.~Erd\H{o}s and R.~L. Graham,
  \href{https://mathweb.ucsd.edu/~ronspubs/80_11_number_theory.pdf}{\emph{Old
  and New Problems and Results in Combinatorial Number Theory}},
  Monogr. Enseign. Math. 28, Geneva, 1980, p.~64.
\bibitem{erdos1988} P.~Erd\H{o}s, \emph{On the irrationality of certain series:
  problems and results}, in A.~Baker (ed.), \emph{New Advances in Transcendence
  Theory}, Cambridge UP, 1988, pp.~102--109,
  doi:\href{https://doi.org/10.1017/CBO9780511897184.009}{10.1017/CBO9780511897184.009}.
\bibitem{kovactao2024} V.~Kova\v{c} and T.~Tao,
  \href{https://arxiv.org/abs/2406.17593}{\emph{On several irrationality
  problems for Ahmes series}}, Acta Math. Hungar. \textbf{175} (2025),
  572--608; preprint arXiv:2406.17593v4.
\bibitem{koizumi2025} J.~Koizumi,
  \href{https://math.colgate.edu/~integers/aa28/aa28.pdf}
  {\emph{Irrationality of the reciprocal sum of doubly exponential sequences}},
  Integers \textbf{26} (2026), Paper No.~A28, 17 pp.,
  doi:\href{https://doi.org/10.5281/zenodo.18714404}{10.5281/zenodo.18714404};
  preprint arXiv:2504.05933v1.
\bibitem{erdosproblems} T.~F. Bloom,
  \href{https://www.erdosproblems.com/243}{\emph{Erd\H{o}s Problem \#243}},
  \texttt{erdosproblems.com/243}, accessed 28 July 2026 (page displays
  ``last edited 21 January 2026'').  The current record labels the problem
  open, cites \texttt{[ErGr80, p.~64]} and \texttt{[Er88c, p.~105]}, and
  explicitly describes its status as the website owner's present assessment
  rather than a literature-completeness guarantee.
\bibitem{formalconjectures243} The Formal Conjectures Authors,
  \href{https://github.com/google-deepmind/formal-conjectures/blob/f776d2f2039351b00737ffcafb9d7d7666e1d9af/FormalConjectures/ErdosProblems/243.lean}
  {\emph{FormalConjectures.ErdosProblems.\texttt{243}}},
  Lean source at commit \texttt{f776d2f}, 2025, accessed 28 July 2026.
  Its $\Q$-valued \texttt{Summable} hypothesis encodes rationality, its
  indexing is zero-based, and its proof ends in \texttt{sorry}.
\end{thebibliography}

\end{document}
